\documentclass[UTF8,a4paper,12pt]{ctexart}

\usepackage[margin=2.35cm]{geometry}
\usepackage{amsmath,amssymb,bm}
\usepackage{booktabs,longtable,tabularx,array,multirow}
\usepackage{xcolor}
\usepackage{enumitem}
\usepackage{listings}
\usepackage{float}
\usepackage{hyperref}
\usepackage{fancyhdr}

\definecolor{reportgreen}{RGB}{36,92,68}
\definecolor{reportgray}{RGB}{245,247,246}
\hypersetup{
  colorlinks=true,
  linkcolor=reportgreen,
  urlcolor=reportgreen,
  citecolor=reportgreen,
  pdfauthor={CUMCM 备赛小组},
  pdftitle={2020--2025 年 CUMCM A/B/C 题学习与论文写作报告}
}
\pagestyle{fancy}
\fancyhf{}
\fancyhead[L]{CUMCM A/B/C 题备赛学习报告}
\fancyhead[R]{2020--2025}
\fancyfoot[C]{\thepage}
\setlength{\headheight}{15pt}
\setlength{\parindent}{2em}
\setlength{\parskip}{0.35em}
\setlist{nosep,leftmargin=2em}
\lstset{
  basicstyle=\ttfamily\small,
  frame=single,
  breaklines=true,
  columns=fullflexible,
  showstringspaces=false
}
\newcommand{\keywords}[1]{\par\noindent\textbf{关键词：}#1}

\title{\bfseries 2020--2025 年全国大学生数学建模竞赛\\A/B/C 题学习与论文写作报告}
\author{队长、组员1、组员2}
\date{2026 年 8 月 1 日}

\begin{document}
\maketitle

\begin{abstract}
本报告面向三人本科生队伍的全国大学生数学建模竞赛备赛。研究范围严格限定为
2020--2025 年 A、B、C 三类题，论文语料为工作区中 59 篇编号 PDF，共 2892 页；
D、E 题不进入统计。报告先按机理建模、综合优化和数据决策三条主线划分个人责任，
随后将思维导图中的算法逐项落实为原理、步骤、代码连接、适用边界和核对产物。论文
分析将编号论文、赛题原文和专家评述分为三类证据：编号论文用于结构和文风统计，
赛题原文用于核对任务，专家评述只用于补充题型背景。全部页面采用文本层提取、
RapidOCR 和 Tesseract 复识相结合的流程，并对公式、表格和小字号图注回看原始栅格。
在此基础上形成十四类章节写作规范、A/B/C 题微观行文差异、创新策略、失分风险与
四周量化训练安排，同时记录双模式论文模板和写作审计 Skill 的实际完成路径。

\keywords{数学建模竞赛；机理建模；组合优化；数据决策；论文写作；证据记录表}
\end{abstract}

\tableofcontents
\newpage

\section{研究边界、证据与处理方法}

\noindent\textbf{语料核验状态（2026-08-01）：}59 篇、2892 页均已完成全文逐页通读，
A/B/C 分别为 20/19/20 篇；每篇均建立十四类章节证据卡，关键公式、表格、小字号图注
和代码页按需回看原始栅格。人工门禁累计覆盖 826/826 个章节类别检查。报告中的模型链、
段落组织和惯用表达来自编号论文；赛题原文只核对任务，专家评述只补充题型背景，不混入
获奖论文文风统计。自动页级统计只用于定位和版面量级，不替代逐页阅读结论。

\subsection{语料边界}

本报告把名称匹配 A/B/C 加编号的 PDF 视为用户提供的国赛一等奖论文语料，不再
外推奖项名单。按年份和题型清点结果见表~\ref{tab:corpus-count}。59 篇论文全部位于
各年份目录根部，页数合计 2892。原始目录 \texttt{2020/} 至 \texttt{2025/} 及思维
导图 \texttt{1.png} 只读使用，不因 OCR、统计或排版而改写。

\begin{table}[htbp]
\centering
\caption{编号论文语料分布}
\label{tab:corpus-count}
\begin{tabular}{crrrr}
\toprule
年份 & A 题 & B 题 & C 题 & 合计\\
\midrule
2020 & 4 & 4 & 5 & 13\\
2021 & 3 & 4 & 4 & 11\\
2022 & 3 & 3 & 2 & 8\\
2023 & 4 & 3 & 3 & 10\\
2024 & 5 & 3 & 4 & 12\\
2025 & 1 & 2 & 2 & 5\\
\midrule
合计 & 20 & 19 & 20 & 59\\
\bottomrule
\end{tabular}
\end{table}

三类材料的用途不可互换。编号论文回答“获奖论文怎样组织和表达”；赛题原文回答
“当年要求解决什么”；专家评述帮助理解命题背景。后两类材料不进入获奖论文的
章节比例、段落密度、图表与公式安排或高频表述统计。2020 年 A/B/C 赛题原文位于
本地 \texttt{2020/problems/} 的 DOCX 文件中，表~\ref{tab:problems} 的任务链据其
正文结构化提取并核对；专家评述仍只用于理解题型背景。

\subsection{页面识别与证据台账}

处理流程以论文为单位建立 JSONL 页面记录。唯一可用文本层的论文直接提取，其余
58 篇先按 120 DPI 渲染并由 RapidOCR 全量识别；文本密集页出现 OCR 中位置信度
低于 0.75、中文正文少于 100 字或标题段落明显断裂时，按 220 DPI 重渲染，并以
Tesseract 的 \texttt{chi\_sim+eng} 复识。每页记录年份、题型、论文编号、页码、
识别方法、置信度、中文字符数、文本框坐标及复识审计信息。

台账把证据细分为章节命中、模型链、写作特征和页面视觉核验四类。公式、表格及
小字号图注不依据 OCR 猜测；这类内容只通过原始栅格的局部裁剪核对。统计脚本必须
先验证 59 个文件和 2892 个连续页码，再计算章节顺序、页数分布、段落密度、图表
候选、公式候选和模型词链。自动页数统计得到全篇页数中位数为 45 页，四分位范围为
37--58.5 页；每篇论文的页面文本行数中位数再取中位数为 38，段落候选数对应
中位数为 6。段落候选
由正文行缩进和垂直间距估计，只用于版面密度比较，不等同于排版软件中的自然段。
共有 2232 页触发并完成 Tesseract 复识，其中 1583 页因质量分更高而被选为最终
文本；最终方法还包括 1281 页 RapidOCR 和 28 页直接文本提取。按同一严格阈值，
1770 页在严格的逐字 OCR 质量阈值下仍标记为未确认；这表示其机器文本不用于逐字引文，
不表示页面未读。59 篇均已逐页人工通读，涉及公式、表格和小字号图注的结论另由原始栅格
或局部裁剪确认。
十四类章节的检出率和页级跨度见表~\ref{tab:section-stats}。

\section{六年真题任务图谱与人员分工}

\subsection{十八道真题的训练定位}

表~\ref{tab:problems} 给出真题与主责能力的对应关系。这里的“主责”表示首先完成
建模和代码基线的人，不表示其他成员不参与复核。每题至少安排一名交叉复核者，
避免模型、数据和论文三条链彼此脱节。

\begin{longtable}{p{1.2cm}p{1.3cm}p{3.3cm}p{6.0cm}p{1.9cm}}
\caption{2020--2025 年 A/B/C 真题训练映射}\label{tab:problems}\\
\toprule
年份 & 题型 & 题目 & 训练任务与模型接口 & 主责\\
\midrule
\endfirsthead
\toprule
年份 & 题型 & 题目 & 训练任务与模型接口 & 主责\\
\midrule
\endhead
2020 & A & 炉温曲线 & 建立焊接区中心温度模型并输出给定工况曲线；求最大过炉速度；在制程界限下分别优化高温面积和曲线对称性。 & 队长\\
2020 & B & 穿越沙漠 & 求已知/当日可知天气下的单人策略，并分别研究已知/未知天气下多人资源、路径、挖矿和购买行为的交互策略。 & 组员1\\
2020 & C & 中小微企业信贷策略 & 量化 123 家有记录企业风险并配置固定总额；评估 302 家无记录企业并制定一亿元策略；考虑突发因素调整。 & 组员2\\
2021 & A & FAST 主动反射面形状调节 & 几何关系、受力/位移约束、连续曲面离散、拟合误差和可实现调节量。 & 队长\\
2021 & B & 乙醇偶合制备 C4 烯烃 & 实验因素建模、反应条件优化、多目标权衡及最优条件解释。 & 组员1，组员2复核\\
2021 & C & 生产企业原材料订购与运输 & 供应商评价、订购计划、运输损耗和多周期约束的联立决策。 & 组员1\\
2022 & A & 波浪能装置最大输出功率设计 & 动力学方程、阻尼与刚度参数、数值积分及平均功率寻优。 & 队长\\
2022 & B & 无人机纯方位无源定位 & 方位几何、编队位置误差、迭代校正与多机协同策略。 & 队长，组员1复核\\
2022 & C & 古代玻璃制品成分分析 & 缺失与异常处理、类别差异、成分关系和未知样品分类。 & 组员2\\
2023 & A & 定日镜场优化设计 & 光学/几何效率、遮挡与截断、场地约束及连续—离散联合优化。 & 队长，组员1复核\\
2023 & B & 多波束测线问题 & 海底几何、覆盖宽度、重叠率、测线方向和总航程优化。 & 组员1，队长复核\\
2023 & C & 蔬菜类商品定价与补货决策 & 销量关联、需求预测、损耗、成本加成与日级补货定价。 & 组员2，组员1复核\\
2024 & A & 板凳龙闹元宵 & 螺线与转向几何、碰撞约束、速度传播和队形可行性。 & 队长\\
2024 & B & 生产过程决策问题 & 零配件抽检、装配、拆解、退换损失和多阶段策略比较。 & 组员1\\
2024 & C & 农作物种植策略 & 地块—季节—作物约束、不确定价格产量和多年度轮作计划。 & 组员1，组员2复核\\
2025 & A & 烟幕干扰弹投放策略 & 三维运动、遮蔽几何、时空覆盖和多平台协同投放。 & 队长，组员1复核\\
2025 & B & 碳化硅外延层厚度确定 & 光谱信号处理、物理测量模型、参数反演和误差分析。 & 队长，组员2复核\\
2025 & C & 无创产前检测时点与异常判定 & 分组时点优化、异常识别、风险约束和模型泛化检验。 & 组员2\\
\bottomrule
\end{longtable}

\subsection{四块学习结构}

队伍采用“三条主责线加一条交叉线”。队长主责机理建模、常微分/偏微分方程、
物理模型、RK4、有限差分、有限元、模型验证和代码集成；组员1主责规划、整数和
多目标优化、智能优化、图论、网络优化、仿真及灵敏度分析；组员2主责数据清洗、
探索、降维、评价决策、回归分类、时间序列和交叉验证。三人共同学习
MATLAB/Python、量纲与误差分析以及 A/B/C 题型选型。共享不等于无人负责：公共
模块仍由对应主责人提供接口，另两人按统一测试样例复核。

\section{理论知识学习笔记}

\subsection{队长：机理、连续模型与数值实现}

\subsubsection{覆盖内容、原理与求解步骤}

队长的模型链从守恒关系开始，而不是从求解器开始。机械系统用牛顿第二定律、
刚体平衡、重力、摩擦和弹簧—阻尼关系建立状态方程；流体问题用质量守恒和
Bernoulli 关系闭合；电路问题用 Ohm 定律和 Kirchhoff 定律组织节点或回路方程；
气体过程核对理想气体状态关系。传热采用 Newton 冷却或热传导方程，波动、扩散
及多孔介质渗流按偏微分方程处理。人口过程用 Logistic 方程和捕食—被捕食模型，
药物过程使用房室代谢方程，传播过程比较 SIR 与 SIS 的恢复机制。

完整步骤固定为：
\begin{enumerate}
  \item 画出对象、边界和能量/质量/动量交换，确定系统尺度与正方向；
  \item 从守恒律或受力关系推导状态方程，逐项给出单位和符号；
  \item 写明初始条件、边界条件、连续性条件、控制量及可行域；
  \item 先做极端情形、量纲和简化模型的解析检查；
  \item 对 ODE 比较显式 Euler 与四阶 Runge--Kutta，对 PDE 按网格选择有限差分或有限元；
  \item 进行步长/网格收敛、守恒残差、回代误差和参数扰动检验；
  \item 把经检验的状态量封装为组员1的目标/约束或组员2的可解释特征。
\end{enumerate}

Euler 法实现简单但全局误差阶数低，只用于基线和步长敏感性对照；RK4 在光滑 ODE
上兼顾精度与实现成本，但对刚性系统不应硬套。有限差分适合规则区域和清晰的
差分格式，可直接检查截断误差与 CFL 条件；有限元更适合复杂几何和变系数边界，
代价是网格、形函数、装配和边界处理更重。解析解存在时先用解析解作单元测试，
不存在时采用网格加密、守恒量和独立实现交叉验证。

\subsubsection{开源代码逻辑、约束与选型边界}

代码流水线按下列接口拆分：
\begin{quote}
\small\ttfamily
geometry -> state -> rhs -> solver\\
-> validation -> export
\end{quote}
MATLAB 侧以 \texttt{ode45}、稀疏矩阵和优化工具箱形成对照实现；Python 侧
以 NumPy/SciPy 的数组、\texttt{solve\_ivp}、稀疏线性代数和网格数据结构实现。
每个模型提供一个可手算的小样例和一个真实规模样例，输出时间网格、状态量、
守恒残差、收敛曲线及下游所需的 CSV。

模型约束至少包括变量域、初边值、守恒关系、几何/运动学限制、参数物理范围和
数值稳定条件。能由规则直接计算的几何关系不使用神经网络；样本驱动关系没有
物理闭合时，不伪造微分方程。ODE 与 PDE 的边界取决于空间耦合是否不可忽略；
有限差分与有限元的边界取决于几何和介质复杂度；RK4 与自适应/刚性求解器的
边界取决于稳定域和尺度差异。

\paragraph{可实现代码清单}
\begin{itemize}
  \item 牛顿运动、弹簧阻尼、刚体平衡、摩擦与重力模块；
  \item Logistic、捕食—被捕食、SIR/SIS、药物房室、Newton 冷却模块；
  \item 热、波、扩散和多孔流方程的 FDM 基线及 FEM 对照；
  \item Euler、RK4、自适应积分接口和统一误差指标；
  \item 网格加密、守恒残差、量纲检查、参数回代和结果导出；
  \item 2022A、2023A、2024A、2025A/B 的最小复现实验。
\end{itemize}

\subsection{组员1：规划、智能优化、图网络与仿真}

\subsubsection{覆盖内容、原理与求解步骤}

规划模型先按变量和约束结构分类：连续线性规划、非线性规划、二次规划、整数
规划、混合整数线性规划和 0--1 规划；再区分单目标与多目标、确定性与随机性。
线性或凸结构优先使用有最优性证书的确定性求解器。整数变量表示启停、选择、
顺序和分配时，必须写出 Big-M 的来源或改用指示约束，不能以取整替代整数规划。
多目标问题先报告 Pareto 关系，再依据题意选择加权、分层或 $\varepsilon$ 约束，
不凭个人偏好给权重。

智能优化覆盖遗传算法、粒子群、模拟退火、蚁群、差分进化和贪心基线，以及
NSGA-II、MOPSO。思维导图中的鲸鱼优化、麻雀搜索、蛾火焰优化作为受控对照：
只有在变量编码、可行性修复和公平预算明确后才进入实验。复杂名称不构成创新，
必须与确定性基线或简单启发式在同一目标、随机种子组和函数评估次数下比较。

图论部分覆盖 Dijkstra、Floyd、SPFA，Prim、Kruskal，最大流、最小费用流、
路径规划、二分图匹配和节点重要性。稀疏非负边单源问题优先 Dijkstra；全源且
规模中等用 Floyd；存在负边时先检查负环，SPFA 只作工程候选而非复杂度保证。
最小生成树解决连通成本，不替代最短路；最大流回答容量，最小费用流同时回答
流量和成本；指派关系优先二分图匹配。

统一求解步骤为：
\begin{enumerate}
  \item 定义决策变量、目标、硬约束、软约束和不确定参数；
  \item 建立可行性基线，并在小规模枚举样例上核对目标值；
  \item 判断线性、凸性、整数性、网络结构和多目标冲突；
  \item 先运行确定性/贪心基线，再决定是否启用元启发式；
  \item 写清编码、初始化、选择更新、约束修复、终止条件和随机种子；
  \item 记录可行率、最优界/间隙、收敛轨迹、运行时间和跨种子分布；
  \item 通过 Monte Carlo 或情景仿真传播参数不确定性，并做敏感性分析。
\end{enumerate}

\subsubsection{开源代码逻辑、优缺点与边界}

代码流水线为：
\begin{quote}
\small\ttfamily
data -> instance -> model -> solver -> repair\\
-> evaluate -> simulate -> report
\end{quote}
Python 用 PuLP/Pyomo 或 SciPy 建立确定性基线，用
NetworkX 验证图算法接口；MATLAB 用 Optimization Toolbox/Global Optimization
Toolbox 建立第二实现。随机算法统一接收种子、预算和停止阈值，输出每代最优值、
可行率及 Pareto 集。仿真模块只通过明确的数据类接收方案，禁止在仿真循环内
偷偷改变约束。

精确规划的优势是可解释、可复算并可给上下界，缺点是非凸或大规模整数问题可能
耗时；启发式的优势是能处理混合和非光滑结构，缺点是没有全局最优保证且对参数
敏感。模型边界由结构而非题型标签决定：可线性化的约束先线性化；不确定性已知
分布可用随机规划或 Monte Carlo，只有区间信息时优先稳健情景；动态图路径与
静态最短路分开处理。

\paragraph{可实现代码清单}
\begin{itemize}
  \item LP、NLP、QP、整数/MILP、0--1 模型生成器及小规模枚举核验；
  \item GA、PSO、SA、ACO、DE、贪心和 NSGA-II/MOPSO 公平比较框架；
  \item WOA、SSA、MFO 的受控插件与相同函数评估预算；
  \item Dijkstra、Floyd、SPFA、Prim、Kruskal、最大流、最小费用流和二分图匹配测试；
  \item Monte Carlo/离散事件仿真、单因素与全局敏感性接口；
  \item 2020B、2021C、2023B、2024B/C 的可行基线。
\end{itemize}

\subsection{组员2：数据处理、评价决策与统计学习}

\subsubsection{数据链、算法原理与完整步骤}

数据工作从数据字典和泄漏边界开始。缺失值按形成机制选择均值/中位数、插值、
KNN 填补或删除；异常值用 $3\sigma$、箱线图和 Isolation Forest 交叉检查，
不能把业务极值自动删除；重复记录先定义主键和时间容差。尺度处理比较
Z-score 标准化与 Min--Max 归一化，类别变量按名义/有序属性选择 one-hot 或
label 编码，连续变量只有在解释或规则需要时才按等宽/等频分箱。样本筛选使用
Monte Carlo 随机抽样时，须固定种子，记录抽样分布并核对覆盖率；它与组员1
的不确定性仿真共享随机数配置，但不把“抽取样本”与“传播不确定性”混为同一任务。

探索阶段采用 Pearson、Spearman、Kendall 相关，分别对应线性、单调秩和序关系；
聚类比较 K-Means、DBSCAN、层次聚类和谱聚类。空间或时间缺口比较 Lagrange、
三次样条和 Kriging 插值。PCA 用于线性方差压缩，t-SNE 主要用于局部结构可视化，
UMAP 兼顾局部与部分全局结构；后两者的二维图不直接充当预测证据。

评价决策链覆盖 AHP、模糊综合评价、模糊 AHP、灰色关联、PCA、因子分析、熵权、
CRITIC、变异系数、TOPSIS、VIKOR、灰色评价和物元可拓，并训练熵权--AHP、
CRITIC--TOPSIS 等组合。主观权重必须给判断矩阵和一致性检验，客观权重必须解释
数据离散或冲突如何影响权重；TOPSIS 的距离、VIKOR 的折中和灰色关联的序列相似
并非同一含义，不能只因都输出排序而互换。

回归覆盖多元线性、Ridge、Lasso、SVR、随机森林、XGBoost 和 CatBoost；分类覆盖
Logistic、朴素 Bayes、SVM、决策树、随机森林、AdaBoost 和 CNN。时间序列覆盖
ARIMA、SARIMA、指数平滑、VAR、GM(1,1)、灰色--Markov、Prophet、XGBoost、
LightGBM、BP、RNN、LSTM、TCN 和 GRU。线性/广义线性模型保留可解释基线，树模型
处理非线性与交互，深度网络仅在样本规模、序列长度和验证方案足以支撑时启用。
思维导图中的“Logistic 预测模型”归入二分类模型，不当作连续响应回归；长期预测是
外推任务，必须报告预测区间、结构漂移风险和超出训练区间的限制。

完整步骤为：
\begin{enumerate}
  \item 建立字段、单位、主键、时间戳、缺失和异常审计表；
  \item 在任何填补、标准化、降维或特征选择前划分训练/验证边界；
  \item 完成探索图表与业务一致性核验，建立简单均值/线性/持续性基线；
  \item 按任务选择评价、回归、分类、聚类或时序模型，并列出近邻方法边界；
  \item 将预处理与模型封装在同一 pipeline，采用 K 折、留出或滚动窗口验证；
  \item 回归报告 MAE、MSE、RMSE、MAPE，分类报告混淆矩阵及类别不平衡指标；
  \item 做残差、校准、特征稳定性、阈值和数据扰动分析，再向决策模块输出。
\end{enumerate}

\subsubsection{开源代码逻辑、约束与选型边界}

代码流水线为：
\begin{quote}
\small\ttfamily
schema -> audit -> split -> preprocess -> baseline\\
-> fit -> validate -> explain -> export
\end{quote}
Python 侧使用 pandas、scikit-learn、
statsmodels 及主流梯度提升接口；MATLAB 侧建立读表、缺失处理、回归分类和时序
对照。所有随机切分记录种子，时间序列禁止随机打乱，分组样本按主体隔离。
评价模型在正向化、无量纲化和权重确定后再计算排序，并做权重扰动和秩相关检验。

K-Means 假定近似球形且需预设簇数，DBSCAN 能识别噪声和任意形状但依赖密度尺度；
层次聚类便于解释树状结构但大样本代价高；谱聚类适合非凸簇但需构图。Lasso 可
稀疏选择，Ridge 对共线性更稳；SVR 适合中小样本但调参和扩展性受限；随机森林
稳健易用但外推弱；梯度提升精度高但要防泄漏和过拟合。ARIMA/SARIMA 要检验
平稳和季节性，VAR 要控制维数和滞后，GM(1,1) 只适合小样本单调趋势，深度时序
模型不能替代滚动窗口基线。

\paragraph{可实现代码清单}
\begin{itemize}
  \item 缺失、异常、重复、尺度、编码、分箱和数据泄漏审计 pipeline；
  \item Pearson/Spearman/Kendall、四类聚类、三类插值和 PCA/t-SNE/UMAP；
  \item AHP/FAHP、熵权、CRITIC、灰色关联、TOPSIS、VIKOR、物元可拓和组合评价；
  \item 线性/Ridge/Lasso、SVR、RF、XGBoost/CatBoost 回归及分类基线；
  \item ARIMA/SARIMA、指数平滑、VAR、GM/灰色--Markov、Prophet 和神经时序对照；
  \item K 折、留出、分组、滚动窗口验证与统一指标报告；
  \item 2020C、2022C、2023C、2025C 的数据基线。
\end{itemize}

\subsection{三人交叉学习与代码集成}

三人共同维护数据契约、配置文件、日志和结果目录。MATLAB 适合快速矩阵推导、
数值求解与官方工具箱复核；Python 适合数据 pipeline、开源模型、自动化测试和
制图。关键结果至少满足“同一语言的独立小样例”或“MATLAB/Python 双实现”之一。
误差分析统一区分测量误差、参数误差、模型结构误差、离散误差和随机算法波动，
不能只给一个综合误差率。

选题时先在两小时内完成任务矩阵：状态量和守恒关系占主导时优先 A 类机理线；
离散资源、路径、排程和策略占主导时优先 B 类优化线；附件数据、预测、评价和
不确定决策占主导时优先 C 类数据线。最终决定同时考虑题意闭合度、可验证性、
团队代码储备和结果文件要求，而不是依据“哪题看起来熟悉”。

\subsection{阶段成果、真实短板与四周计划}

\subsubsection{当前可取之处}

本轮已经形成全算法责任表、18 题训练接口、统一代码分层、模型选型边界和论文
证据流水线。“算法数量多”不在本轮评价范围内；可取之处仍是每条主责线都要求基线、约束、
独立检验和下游接口。复杂模型必须与简单模型公平比较，模型、代码和论文中的
变量名可以追到相应结果文件。

\subsubsection{当前真实短板}

成员未提供个人既有学习记录，因而本报告不虚构“已经熟练掌握”的经历。当前能
确认的是计划与分析产物，不能确认的是三人在限时条件下的独立实现能力。具体
风险包括：队长尚缺复杂边界 PDE/FEM 的限时网格调试证据；组员1尚缺大规模
MILP 的界、间隙与不可行诊断记录，也缺元启发式跨随机种子统计；组员2尚缺严格
时间/分组隔离的数据验证记录，深度时序和客观评价权重容易停留在调用层面。
全队尚未完成一场 72 小时全流程演练，论文数字、图表和附件代码的一致性也没有
实战验收。这些短板按“尚无证据”表述，不把未知情况写成个人缺陷。

\subsubsection{2026 年 7 月 29 日至 8 月 25 日量化训练}

\begin{longtable}{p{1.4cm}p{1.55cm}p{3.1cm}p{4.6cm}p{3.4cm}}
\caption{四周训练计划与验收标准}\label{tab:four-weeks}\\
\toprule
时间 & 负责人 & 训练任务 & 产物 & 验收标准\\
\midrule
\endfirsthead
\toprule
时间 & 负责人 & 训练任务 & 产物 & 验收标准\\
\midrule
\endhead
7/29--8/4 & 队长 & 完成 2022A 动力学最小复现；比较 Euler、RK4 与自适应积分。 & 方程说明、两套代码、步长收敛图、守恒与回代误差表。 & 小样例误差随步长按预期下降；关键结果可复算；接口测试全部通过。\\
7/29--8/4 & 组员1 & 完成 2024B 的枚举基线与 0--1/MILP 表达。 & 决策树、模型文件、枚举核验、不可行情形清单。 & 小规模最优值一致；每条约束有题意来源；报告最优界或间隙。\\
7/29--8/4 & 组员2 & 完成 2022C 数据审计、可视化与分类基线。 & 数据字典、清洗审计、pipeline、验证报告。 & 预处理只在训练集拟合；异常处理可追溯；至少三类基线同口径比较。\\
8/5--8/11 & 队长 & 完成热/扩散方程 FDM 与复杂几何 FEM 对照。 & 网格文件、边界测试、网格收敛表。 & 至少三档网格；边界和量纲单测通过；说明 FDM/FEM 选型边界。\\
8/5--8/11 & 组员1 & 完成 2023B 图/几何基线及 GA、PSO 公平对照。 & 基线求解器、随机算法报告、可行性修复测试。 & 不少于 20 个种子；预算相同；报告中位数、离散度、时间和可行率。\\
8/5--8/11 & 组员2 & 完成 2023C 时序滚动验证和补货特征接口。 & 持续性/季节基线、ARIMA、树模型、滚动评估。 & 无未来信息泄漏；至少四个误差指标；失败日期有原因分析。\\
8/12--8/18 & 全队 & 选择一题进行 36 小时压缩演练，贯通模型、代码、图表、论文和附件。 & 可编译论文、运行入口、结果清单、复盘表。 & 所有正文数字可定位到输出；无未定义引用；第三人按说明复现关键表。\\
8/19--8/25 & 全队 & 进行 72 小时 A/B/C 三选一全流程模拟并交叉审稿。 & 最终论文、源代码、附件、审计 JSON、个人复盘。 & 12 小时内定题；24 小时出基线；48 小时主结果冻结；终稿通过结构、引用、检验和复现门。\\
\bottomrule
\end{longtable}

\section{59 篇编号论文的结构与写作分析}

59 篇论文的自动页级候选给出篇幅与标题统计；59 张人工证据卡确认逐篇模型链、
章节写法、结果解释、检验和创新结论。自动候选只标出页级位置，人工卡用于判断
段落功能和真实行文动作，两类证据不互相替代。

\subsection{宏观框架与版面组织}

编号论文并不追求固定的小标题数量，而是把题目信息、建模依据、求解过程、结果解释和误差核查逐项写清：
\[
\begin{aligned}
\text{任务界定}&\longrightarrow\text{假设与变量}
\longrightarrow\text{模型及约束}\longrightarrow\text{求解实现},\\
&\longrightarrow\text{结果解释}
\longrightarrow\text{独立检验}\longrightarrow\text{边界与复现材料}.
\end{aligned}
\]
章节标题可以因题目合并，但上述功能不能被删去。篇幅判断使用全量论文的实测
页级位置和章节跨度，不把经验百分比写成硬规则。全篇总页数参考中位数
45 页和四分位范围 37--58.5 页；不同章节的检出率、首次出现位置和跨度汇总于
表~\ref{tab:section-stats}。检测率低的章节只能形成
定性建议，不能据缺失 OCR 推断获奖论文没有相应内容。

\begin{longtable}{p{1.8cm}p{2.5cm}p{3.0cm}p{2.8cm}p{3.0cm}}
\caption{十四类章节的全语料检出与页级位置统计}\label{tab:section-stats}\\
\toprule
章节 & 检出论文 & 首次位置比例中位数 [Q1,Q3] & 跨度页数中位数 [Q1,Q3] & 跨度比例中位数 [Q1,Q3]\\
\midrule
\endfirsthead
\toprule
章节 & 检出论文 & 首次位置比例中位数 [Q1,Q3] & 跨度页数中位数 [Q1,Q3] & 跨度比例中位数 [Q1,Q3]\\
\midrule
\endhead
摘要 & 58/59 (98.3\%) & 0.0215 [0.017, 0.0268] & 1 [1, 1] & 0.0215 [0.017, 0.0268]\\
问题重述 & 57/59 (96.6\%) & 0.0444 [0.0345, 0.0541] & 1 [1, 1] & 0.0222 [0.0172, 0.0286]\\
问题分析 & 57/59 (96.6\%) & 0.0526 [0.0417, 0.0732] & 1 [1, 1] & 0.0244 [0.0172, 0.0357]\\
模型假设 & 48/59 (81.4\%) & 0.0735 [0.0562, 0.106] & 1 [1, 1] & 0.0225 [0.0172, 0.0286]\\
符号说明 & 58/59 (98.3\%) & 0.0808 [0.0661, 0.1143] & 1 [1, 3.75] & 0.029 [0.0205, 0.0749]\\
模型建立 & 58/59 (98.3\%) & 0.0825 [0.0489, 0.1329] & 2 [1, 6] & 0.045 [0.0245, 0.1681]\\
模型求解 & 55/59 (93.2\%) & 0.1077 [0.0286, 0.2319] & 2 [1, 6] & 0.0385 [0.0206, 0.1562]\\
结果分析 & 47/59 (79.7\%) & 0.1795 [0.0639, 0.3107] & 7 [1, 12] & 0.1282 [0.0278, 0.25]\\
模型检验 & 22/59 (37.3\%) & 0.2294 [0.0564, 0.4821] & 1.5 [1, 7] & 0.0335 [0.0249, 0.1889]\\
灵敏度分析 & 25/59 (42.4\%) & 0.0357 [0.02, 0.3929] & 1 [1, 2] & 0.0256 [0.0179, 0.0385]\\
模型评价 & 51/59 (86.4\%) & 0.5111 [0.396, 0.6186] & 1 [1, 1] & 0.0227 [0.0172, 0.0294]\\
改进方案 & 19/59 (32.2\%) & 0.4667 [0.3707, 0.5928] & 1 [1, 1] & 0.0263 [0.0194, 0.0345]\\
参考文献 & 57/59 (96.6\%) & 0.5424 [0.4151, 0.6667] & 1 [1, 11] & 0.0333 [0.0222, 0.2941]\\
附录 & 48/59 (81.4\%) & 0.5574 [0.4427, 0.7079] & 16 [6, 28] & 0.3779 [0.1865, 0.4926]\\
\bottomrule
\end{longtable}

表中首次位置按页码除以论文总页数计算。灵敏度分析的首次位置中位数异常靠前，
抽查发现目录页和摘要中的同名词会造成提前命中，故该项只用于检索，不能解释为
正文真的在前 3.57\% 篇幅内展开灵敏度分析。参考文献与附录跨度也可能因标题合并、
附件页连续而偏大；模板只把高检出章节的跨度用于量化提示。

\subsection{逐篇证据摘要（59/59）}

表~\ref{tab:verified-paper-cases} 压缩全部已经全文通读并核验关键栅格的五十九篇论文；
每行只保留模型链、论证次序和可迁移的人类作者表达。需要限制结论强度的事实仍记在
单篇证据卡中，不在汇总表内追错，也不把个别现象解释为 59 篇中的出现频率。
每篇完整的十四类章节证据、页码定位、结果摘录和撤回记录保存在 Skill 的
\path{references/paper-cards/}，报告随人工门禁进度逐篇增量更新。

\begin{longtable}{p{1.4cm}p{5.2cm}p{7.0cm}}
\caption{五十九篇已核论文的模型链与可迁移写法}\label{tab:verified-paper-cases}\\
\toprule
论文 & 模型链 & 可迁移写法\\
\midrule
\endfirsthead
\toprule
论文 & 模型链 & 可迁移写法\\
\midrule
\endhead
A070 & 热传导 PDE 候选、有限差分与识参、集总 ODE 降阶、制程约束优化。 &
先保留高保真候选，再用误差和计算成本说明降阶；摘要与各问结果逐项闭合。\\
A147 & 一维热传导 PDE、有限差分、四参数拟合、爬山热启动的限时队列搜索。 &
以已得可行解热启动后续搜索，并把模型、搜索和制程读数串成一条实现链。\\
A195 & 环境温度插值、热传导 PDE、Crank--Nicolson 型离散、标定、粗细搜索、7 次 GA、分层种群接力。 &
把边界写入离散系统；利用单调性缩小范围；跨层保留末代种群并报告多次随机运行。\\
A212 & 线性环境温度场、集总指数升温与牛顿冷却、残差驱动分段修正、参数扫描、粗细网格、退火与加权。 &
先建低阶可计算基线，再按残差时段提出机制修正；逐问报告制程硬指标，并提供从状态更新到目标函数的附录链。\\
A028 & 理想抛物面参数化、三维旋转、节点/球面弧中点 RMS 微调、面板蒙特卡洛接收比与球面基线。 &
保留 300.4 m 实际半径和球面面板，按“基准曲面--方向变换--离散/连续几何误差--任务指标”逐层组织，并同时报告微调前后 RMS 和原始球面基线。\\
A115 & 理想抛物面可行性试算、变步长极小极大搜索、偏轴截面样条映射、面积比接收估计与单参数扫描。 &
先报告最大邻距变化率与通过阈值的节点比例；把偏轴几何的“提案、实现、求解”状态分开；用基准接收率和扰动曲线暴露代理模型的待检验环节。\\
A217 & 轴对称抛物面机理、焦距一维驻点、坐标旋转、带结构约束的高维最小二乘、三角面板射线追迹与蒙特卡洛积分。 &
先抓住焦距这一低维自由度；第二问通过坐标旋转复用径向关系，第三问再把相交、法向、反射、目标平面和圆盘命中逐层闭合。摘要一问一段，成对图分别解释形面、执行器、误差和基准。\\
A001 & 波浪激励下浮子/振子垂荡和纵摇受力平衡、一阶状态 RK4、稳定窗口平均功率、阻尼粗搜到 GA 精求、双阻尼与周期检验。 &
先从坐标和受力写出状态，公式后紧跟物理解释；由前期不稳定观测确定 100--180 s 稳定窗口，先粗略遍历再用 GA 求精。\\
A022 & 波浪能装置平衡水平、振动方程、恒定/幂律阻尼、垂荡--纵摇耦合、稳态积分、阻尼粗搜与局部精化。 &
先用平衡高度排除浸没边界，再由第一问数值结果授权复杂耦合；恒定阻尼先给复数解再用 trapz，阻尼和功率按粗搜、局部加密和排序收束。\\
A171 & 浮子/振子平衡与分段浮力、常系数拉普拉斯解、幂律阻尼 RK4、力电比拟、稳定窗口功率、垂荡--纵摇八状态模型与双阻尼优化。 &
先写平衡位置和浸没边界，再按方程性质切换解析解与 RK4；功率先定义稳定窗口，结果表沿用同一组时刻，检验把瞬态误差和稳态误差分开判断。\\
A0127 & 双坐标反射几何、五类光学效率、塔影中心判据、镜面无效点与锥形射线追迹、分区同心圆布场、单位面积输出优化、二分/变步长搜索和灵敏度。 &
先把三问写成正向计算、反问题和放宽约束，再由题面已知量隔离五类关键效率；工程简化按计算负担、数量占优、中心对称与统计补偿给依据，云图直接转成塔址和布场方向，搜索按计算预算分粗细两级。\\
A0165 & 矩形点判定、太阳位置与镜面法向、六近邻遮挡栅格、锥形光斑截断、五类效率、空间图选向、固定变量降维、三分搜索和按圈高度蒙特卡洛。 &
摘要按光传播过程分段；把矩形判定化为两个三角形；先用 $2\times2$ 快速核并报告不足 $2\%$ 的偏差，再固定其余变量读单峰响应；空间图直接触发塔向南移动，问题三则冻结宽度、塔位和布场，只开放各圈高度。\\
A0175 & 三坐标系与锥形光束、塔影/入射阴影/反射遮挡、20 m 邻接矩阵、集热器接收、蜂窝布局、整体旋转、两级镜面抽样和外圈高度调整。 &
把程序短路写成“若塔影成立则停止，否则继续检查邻镜”；先给全对全计算量再引出邻接候选；用六档步长的增幅趋缓选择 1.5 m；以固定基线试算后冻结旋转角和高度；明确写出 $10\%\times10$ 粗评与 $30\%\times5$ 细评，并在超过 60 MW 后继续讨论删去高遮挡镜面。\\
A092 & 太阳与镜面坐标变换、入射/反射遮挡投影、重叠区域离散并集、万次锥束蒙特卡洛、空间方向试算、统一规格 GA 与按圈变步长优化。 &
先逐项排除可直接计算量，只把唯一不能直接求得的截断效率交给随机估计；由环序号曲线提出北侧高效猜想，再选代表时刻二维图复核；东西方向关于原点近似对称且上午、下午变化抵消后冻结横向坐标；第三问由前问同心圆结构和各向同性把逐镜变量压为逐圈变量。\\
A016 & 等距螺线弧长、相邻把手定长与速度递推、矩形投影碰撞、首次事件搜索、最小螺距回退、两圆弧调头几何、四区域状态矩阵和速度齐次缩放。 &
先把全局运动压成相邻把手递推；多根由物理排列顺序和局部小角度裁决；由轨迹继承把首发碰撞对象缩到龙头/第二板凳的少量顶点；搜索保留越界、回退、折半和误差界；题设的缩短目标若被几何不变量锁死，先证明不可实现，再继续指定方案。\\
B078 & 动作剪枝、资源/现金状态递推、已知天气图搜索、未知天气期望决策、候选路线支付矩阵与多人动态定义。 &
按已知/未知天气、单人/多人拆分信息结构；将移动、停留、挖矿和补给写成可对照代码的状态转移，并逐问给出结果表和附录入口。\\
B108 & 游戏引擎与图预处理、动态规划状态/价值压缩、1024 种未来天气枚举、统计规则、蒙特卡洛评估、天气等价换算、候选路径与局部日博弈。 &
把现金作为同资源状态下的价值并做支配压缩；分开报告失败率、成功条件均值和失败记零的无条件均值；由支付矩阵逐项核对最佳反应与均衡。\\
B125 & Floyd 图约简、资源状态动态规划、天气马尔可夫链、完整天气情景 DP、人工势场/拥挤收益与静态支配判定。 &
把功能节点最短路和资源状态转移分层；将完整未来情景 DP 限定为离线 oracle，再以同一非预见策略在留出情景上评估遗憾与库存风险。\\
B175 & 关键节点最短路、确定天气多阶段策略、位置状态 MDP、天气马尔可夫链、单路径模拟与多人协作目标。 &
逐层区分确定天气、未知天气和多人信息结构；把效率与公平明确写成两个目标，并由附录反查报告值的计算来源。\\
B007 & 受控条件分组、Pearson/二项式拟合、Q 型聚类、多元回归与偏最小二乘比较、交互收率优化、均匀设计。 &
四问按“认识已有数据--解释变量作用--求最优条件--设计下一轮实验”递进；先写旧模型遗漏，再加入直接收率回归和温度交互项；实验预算、安全修正、备选方法与舍弃理由均明确出现。\\
B026 & 三次样条统一温度水平；散点与转化率上界引出 Logistic；一次/二次曲线与分段函数描述选择性和时间变化；双因素方差分析、多重比较；五元/三元二次回归及分情景优化；五次针对缺口的新增实验。 &
先写数据整理的麻烦，再用散点形态和机理边界共同决定模型；每次实验都交代缺口、保持条件、改变因素、观测量和裁决用途；无温度限制与受限结果成对呈现。\\
B050 & 42 张散点分类与线性/指数、三/四次多项式裁决；反应时间过程；六变量递进回归；BP 收率优化；加权欧氏距离与五次分角色实验。 &
用样本点数解释为何舍弃 $R^2=1$ 的四次插值；每轮回归以新旧指标推进；先写规划方案的具体不足再切换 BP；实验先分探索与确认。\\
B160 & 每组一元二次回归、控制变量双响应比较、混合水平正交极差、七变量二次交互回归、逐步回归、收率回归、\texttt{quadprog} 与五次分角色补测。 &
曲线选型同时写明线性不足与高阶边界；控制变量写全保持量、改变量和组别；诊断点名共线性后再修模；回归按相对、绝对和方差分析三层解释。\\
B030 & 极坐标下三类九形态、旋转对称压缩、六个正弦定理定位分支、误差遍历、发射机编号识别、三步队形调整与锥形编队圆形化。 &
先穷尽情形再由可见对称关系压缩；算法段写明固定量、遍历量、范围和输出；旧模型复用先交代授权，结果按读数、原因和整体对照推进。\\
B035 & 极坐标正弦模型、定弦定角两圆轨迹、分层编号解码、短弧与角度区间穷举、三/四机定位、$L_2$ 前视启发式搜索、退化处理与随机误差仿真。 &
先写分类方程的可见负担再换几何轨迹；证明按数值例、区域交叠、端点界和区间表推进；算法比较允许局部让步；数值极限与实际停止分开说明。\\
B086 & 极坐标正弦定位、解析双解唯一化、正九边形角集合判号、径向预处理、两类局部平方角误差搜索、三引理锥形调整、平面拟合与高斯扰动。 &
把局部可见信息与事后全局评价分栏；证明误差界后继续解释搜索与通信收益；方案按收敛、稳态误差和资源代理裁决；几何引理由控制量映射到具体机群。\\
B226 & 正弦定理覆盖几何、三维线面夹角、同长度覆盖面积证明、边界锚定贪心、模拟退火对照、等深线分区、最小二乘局部平面、PSO 与残差驱动再分区。 &
先把三维问题归约为水深和线面夹角两个缺失量；命名平面后推导向量；规范只作动机并另给命题证明；贪心写清东边界首线、1 m 步长和最低可行重叠率；拟合残差只触发失败区域再分；按异常区域分别解释线数与漏测。\\
B311 & 两种重叠定义与 CAD 裁决、任意方位线面几何、长尺度累积、边界递推和真实海域最不利位置覆盖。 &
定义有歧义时保留两套解释，以同一数值例和可观察图形裁决；跨问只改水深与夹角输入；小坡度先放到 2 海里尺度计算；题内定量判断先于规范印证；递推写到对边闭合并舍弃额外候选。\\
B477 & 投影重叠率修正、二维关系升维、平行测线递推、随机森林地形接口、局部梯度搜索与等深线分区修复。 &
先用共同输入下的反直觉排序迫使定义修正；航向固定时先证明坡度唯一再升维；离散加密按精度、失效、误差累积和耗时列出负担；初始搜索暴露起点依赖与漏测后，只修浅缓失败区；主文示意、附件坐标和分区总量各司其职。\\
C050 & 数据侧写、负价退货剔除、相关与四级分群、成本加成回归、半年窗口季节预测、品类/单品两级模拟退火、语义特征灰色关联和外部数据分流。 &
每张侧写图说明后续用途；代表单品先给排序量与阈值；回归斜率译成每增加 $1\,\mathrm{kg}$ 的定价变化；丰富度、节气和节日指标写明代理语义；附件内证据与待采数据分开，打折频率按竞争机制、追加证据和分情形动作解释。\\
C126 & 日销售透视、同期平均、T/F 相关矩阵、正态/对数正态 Bayesian 销量模型与 MCMC 诊断、二阶 VAR、两路单品筛选、利润优化和两步 LINGO 补货定价。 &
先以不正态、大量零值、季节性和计算负担授权数据变换；由隔日难售的经营事实确定两天记忆和 VAR(2)；把负预测译成余货、折价与不补货；正文保留代表规则和关键诊断，完整 BUGS、R、EViews、LINGO 入口转入支撑材料。\\
C228 & 数据连接与侧写、时间/销售/关系三路聚焦、偏相关与 FP-Growth、双对数价格弹性、候选预测比较、LSTM、弹性修正、GBest-PSO、VIKOR 与 NSGA-II 多目标单品决策。 &
开放题先说明未指定探究方向，再把观察面收窄为互补问题；统计周期后接节日、作息和供给季节机制；ARIMA、回归和灰色预测效果不佳后才改用 LSTM；缺失、短窗无趋势和随机波动又触发加权平均；变量变为 $0/1$ 与连续混合后再说明 NSGA-II 和混合编码。\\
C006 & 六指标熵权--TOPSIS、二元双目标筛选、GP 供货误差、动态库存/成本规划、逐周订购 GA、规则式转运、A/C 原料结构调整与扩产。 &
把潜变量代理、拒绝历史均值、状态传递、宏观/微观解释、有限改进和模块沿用写出完整理由。\\
C085 & 四指标熵权--TOPSIS、稳定供货筛样、动态权重贪心/动态规划、分阶段订购与转运、偏差变量规划、二分产能上界与窗口灵敏度。 &
先把指标方向接到业务能力；由固定权重的单周缺货风险改权重，由连续淡季定两周窗口；原始成本不可比时改用成本/产能；求解器按预拆、背包和残项补救接力。\\
C169 & 八指标投影寻踪与加速遗传优化、逐周 0--1/线性规划、MOPSO、多阶段订购转运和相邻周恢复约束的产能规划。 &
把业务后果的不对称性先写成函数条件；观察到规律后说明为何不纳模；用历史同周期位置确定初值；先判时间记忆再拆周；按任务语义对 24 个周结果取最大或最小。\\
C283 & 九指标熵权与 AHP 有限修正、等效产能库存、随机供货/损耗双目标、双层贪心、候选订货式比较、定向材料替换和扩围产能。 &
指标写到具体数据操作；客观权重出现业务异常后才修正；用 LP/DP 边界说明贪心选型；停止规则写明继续下单的代价；随机运行只选高频较优代表；目标逆转时同步改变候选范围。\\
C155 & 成分闭合与中心化对数比、Pearson/Yates 前提分流、风化阶段 Q 聚类、逐成分二次回归与逆变换、分层 PbO 决策树、R 型变量聚类、Q 型样本聚类和灰色关联。 &
先按缺失字段语义决定剔除、置零或不填补；检验前先看期望计数；用前问结论命名聚类并反向核对；按 $R^2$ 分流预测与不变策略；变量聚类和样本聚类分别承担降相关与分组。\\
C229 & 条件化颜色填补、近似零值乘法替换、闭合与 clr、Spearman/卡方、Aitchison 均值、Dirichlet 回归、决策树--PLS-DA 复核、WSS--K-means、交叉验证 PLS 和配对 Wilcoxon。 &
同一字段先按候选样品参考质量分流热卡与众数；主导成分遮蔽时保留全图并补画去除图；先用 PbO 阈值给直观规则，再以 PLS-DA/VIP 复核；没有亚类标签便说明无监督身份，聚类不清时只称粗略三分；潜变量由 RMSEP 与解释率共同裁决。\\
C109 & 发票清洗与年度补全、七项企业特征、XGBoost 违约风险、评级--风险级联、信贷规划及宏微观疫情调整。 &
把交易数据、可解释特征、风险概率和逐企业额度/利率串成决策链；摘要逐问报硬结果，长表与代码提供追踪入口。\\
C142 & 发票清洗、七指标、PCA 安全指数、违约/流失函数、遗传算法、两级 BP、行业/事件冲击矩阵与重新优化。 &
将有信用记录与无记录企业接入同一策略链，并对贷款总额、违约函数方差和冲击比例做三组敏感性；附录覆盖指标、预测和优化实现。\\
C170 & 发票清洗、十一指标、熵权/TOPSIS 风险、经验阈值、RAROC 式额度规划、Logit、BP 与三类事件因子。 &
把企业画像、风险标签补全、放贷筛选、额度分配和事件重求解按三问串联；附录公开指标汇总、权重、利率分组和行业分类代码。\\
C227 & 四层 20 特征、五分类器 SoftVoting、评级风险聚合、多目标信贷规划、XGBoost 评级、层次聚类与疫情风险乘数。 &
按“企业画像--违约概率--信誉评级--额度/利率--冲击重算”连接三问；同一金额口径下组合收益、违约损失和客户流失，并公开预测与规划代码片段。\\
C305 & 六指标 Logistic 守约概率、等级均值最近距离、单/多目标贷款规划与外部冲击折减。 &
将企业特征、守约概率、等级、额度、利率和冲击后重规划串成三问，并公开清洗、分类和规划代码。\\
A053 & 螺线弧长、定长递推、首次碰撞反证、矩形分离轴、变步长螺距搜索、相切圆弧、四路径状态与 PSO 限速。 &
多根筛选写成“候选根--构件次序--方向条件--最近可行根”；由首次事件对象切换解释参数微变后的跃变，并以“龙身不能倒退”的现实后果推出半径下界。\\
A163 & 螺线速度递推、碰撞域缩减、四角点判定、步长搜索、圆弧不变量、四路径状态与速度比。 &
先证明几何自由度，再从全局轨迹图缩到局部搜索；算法段按状态类别、运动方向和推进步长交代递推。\\
A178 & 代理圆预筛、四角点精判、螺距二分、圆弧不变量、四时段递推、曲率猜想、多速度核查与局部二分。 &
代理模型只负责夹住区间，精确几何负责最终裁决；先提出机制猜想，再跨条件核验，不把单幅图形直接写成规律。\\
A242 & 首次碰撞缩域、点线距离、螺距 PSO、盘入盘出路径、齐次速度传递与粗搜--三分求界。 &
优化区间由物理下界和先验可行上界共同给出；仿真、几何判定和速度约束各承担不同验证任务。\\
A196 & 视线球面布尔积分、连通区间证明、凸截面、圆周离散、二分、误差界、多初值、差分进化与整数分配。 &
先证明“只进出一次”再授权二分；理论误差界与最终输出精度分开，近似方案只提供初值而不冒充最终解。\\
B159 & 检验假设反向、功效与成本定样本、16 种检测方案、互斥成本核算、动态规划与 Beta 后验积分。 &
先由企业对风险的态度确定原假设与备择假设；复杂度增加时明确说明为何换模型，后问只替换不确定性输入而保留决策骨架。\\
B195 & 无实测样本下的模拟器、$(n,c)$ 抽检、跨轮状态、重复步骤压缩、长短流程权衡与后验区间。 &
先如实交代数据缺口，再把时间和费用核算闭合；公开“原步骤重复，故由 27 步压至 16 步”的修正过程，保留长期与短期偏好的真实取舍。\\
B196 & 节点与反馈计数、两级成本、蒙特卡洛/蚁群、遗传算法和节点扰动分析。 &
先用结构规模授权搜索方法切换；由高频节点落到具体执行次序，再按“扰动--变化节点--成本原因”解释敏感性。\\
B060 & 复折射率与 Fresnel 方程、局部/全曲线反演、网格与 Gauss--Newton、Airy 公式、判据和误差传播。 &
预处理的两个目的分别交代；局部路线与全曲线路线按证据范围分工，中间谱差必须重新传入厚度模型后才解释最终影响。\\
B157 & 双角度共享物理量与独立标定、内层闭式解、外层非线性拟合、残差触发的 $C^1$ 分段与 Airy 修正。 &
先说明频率干扰会怎样伤害下游反演，再决定是否去除；主结果基本不变而局部残差下降时，只把修正限定为局部作用。\\
C038 & 七年种植规划、十四类约束、DEGA、1000 情景与 95\% CVaR、分布/相关检验和联合扰动。 &
图表观察直接进入变量和约束；风险下降与利润牺牲同时报告，先说明样本量和检验前提，再给方法；联合不利情景与单因素扰动分开写。\\
C063 & 基准规划、随机状态、典型候选、共同环境复评、收益矩阵、风险准则与替代互补关系。 &
典型情景用于解释，共同随机环境用于裁决；先写最坏后果再给极大极小选择，表后按幅度与优先级解释方案差异。\\
C094 & 决策量 $X/Y/Z$、管理参数 $p/q$、Gurobi/贪心、共同环境复评和替代互补扩展。 &
把管理口号落实为 $p/q$ 参数；候选方案先生成，再放入同一环境比较；只有需求关系和目标响应都成立时才引入替代性。\\
C234 & 参数索引、十四类约束、染色体可行生成与修复、灾害链和供需联动。 &
先用数据表决定索引，再把地块、季次与作物三类邻接关系写进编码；按收敛、约束形状和极端基线三条证据说明结果可信度。\\
C023 & Logit 混合效应、阈值概率、动态规划分组与时点、连续/阶梯损失、稳健空值与 FDR、三级风险判别。 &
先把安全、及时和可执行三种损失说清，再区分概率校准与决策效果；“结果可疑”在流程中被翻译为复测，而不是一句模糊警告。\\
C132 & 重复测量混合模型、个体 GPR 与三级回退、BMI 遗传分组、MLP/LSTM/DeepHit、生存风险和 OOF 专家融合。 &
先把多行检测转换为孕妇级时点；模型分工按首次就诊特征、时序变化和删失风险展开，最终同时报告总体表现、少数类召回和样本数。\\
\bottomrule
\end{longtable}

段落组织以“一个判断配一组证据”为单位。建模段通常先说明关系来源，再定义变量
并列方程，公式后解释每一项的作用；求解段按输入、初始化、迭代、停止和输出展开；
结果段按读数、比较、原因和决策含义展开。图表应在首次讨论附近出现，正文必须
引用并解释，完整中间矩阵和长代码移至附录。公式、表格和图注的数量只能作为
页面候选统计，最终含义以原始页面裁剪核验为准。每篇图表标题候选数的中位数及
四分位范围为 17 [14, 23]，公式行候选对应为 420 [228.5, 511.5]；两类候选页在
全文相对位置的中位数分别为 29.8\% 和 56.9\%。这些计数用于判断版面节奏，不表示
精确的图表或公式数量。

\subsection{A/B/C 三类论文的微观文风}

分题型比较采用各类论文内部的章节、连接词和模型组论文检出数，不将某个模型词
出现次数直接解释为模型质量。A 类总页数中位数和四分位范围为 44.5 [36.5, 58]
页，模型组检出前三项为机理与数值 15 篇、规划与智能优化 13 篇、时序预测 3 篇；
B 类为 42 [34, 55] 页，前三项为机理与数值 12 篇、规划与智能优化 10 篇、统计
学习 6 篇；C 类为 52 [43.25, 60.25] 页，前三项为时序预测 17 篇、规划与智能
优化 17 篇、统计学习 13 篇。模型词可能来自摘要、比较方法或参考文献，这里只报告
“论文中检出”，不把它直接解释为论文主模型。

\subsubsection{A 类：先交代机制，再讨论数值}

A 类行文从研究对象、坐标、尺度和守恒关系进入。高信息密度表达包括“以……为
研究对象，取……方向为正”“由质量/能量/动量守恒可得”“在边界……处满足……”
“将区域离散为……，当网格由……加密至……时，目标量变化……”。推导篇幅不必刻意拉长，
假设、方程、初边值和数值格式仍要彼此闭合。结果解释需回到物理量、数量级和机制；
一条平滑曲线本身不足以说明模型正确。

需要避开的写法是先报“使用微分方程模型”，随后直接给软件输出；或者给出方程
却没有边界、单位、稳定条件和网格收敛。图形以结构示意、网格、状态曲线和残差/
收敛图为主，颜色只服务于变量区分。

A070、A147、A195、A212、A028、A115、A217、A001、A022、A171、A0127、A0165、
A0175、A092、A016、A053、A163、A178、A242 与 A196 共二十篇 A 题均已通过人工门禁。
这些论文最值得学习的共同点，是把“为什么可以这样算”留在正文中：守恒关系授权方程，
边界条件限定数值格式，几何次序筛去无意义的根，现实动作排除数学上可行但物理上不可能的
分支。作者不急于给算法名，而是先写观察、冲突或单调关系，再自然地落到求解方法。

A028 与 A115 的可迁移价值在于把理想几何、工程约束和执行方案分层表达。前者先以坐标和
面板关系构造几何候选，再用索长、邻距和受力兼容性判断能否执行；后者把节点逼近、连续
工作面和接收效率分成三个口径，提醒写作者不要用一个代理量代替完整物理链。仿写时应保留
“先得到候选，再检查……”这一顺序，并在结果段同时交代几何读数、约束余量和工程含义。

A053、A163、A178 与 A242 围绕同一类板凳龙任务，却给出了四种很像人类研究的推进方式。
A053 用完整反证把首次碰撞逐节前推到龙头；A163 先证明自由度，再由全局轨迹图缩小局部
搜索域；A178 让代理圆只负责夹住区间，最终仍回到四角点精判；A242 先给物理下界和经验
可行上界，再在有来由的区间内优化。四篇都把“缩域依据”写在“搜索方法”之前，避免读者
看到一个凭空出现的上下界。

其中 A053 的结果解释尤其值得直接转化为写作习惯：参数只变化 $0.01\,\mathrm{cm}$ 时，先查
首次触发事件是否从一个构件切换到另一个构件，再解释输出为何跃变；某个终点方案看似可行，
则把它放回全过程寻找提前碰撞，以反例裁决简化方案；半径继续减小时若要求后续构件倒退，
就由这一现实后果推出数学下界。A196 则先证明目标区间连通，才使用二分；理论误差、输出
精度和多初值结果分开报告。它们共同说明，模型可信度常来自一两句关键判断，而不是更长的
算法介绍。

A217 展示了机理几何题更值得迁移的一面：先由轴对称抛物面的截面关系把焦距问题
压缩成一维驻点，再沿用同一坐标框架处理面板姿态和遮挡，最后以三角面板射线追迹
配合蒙特卡洛积分估计接收效率。正文不靠长串套话衔接，而用“注意到……，故……”
交代关键观察，用“考虑到……，本文将……”给出取舍；公式之后马上说明几何量的
含义，成对图形则以“图左/图右”分别解释基准状态与调整状态。这种写法把观察、
选择、推导和图证压在同一条论证线上，比先报算法名称再补理由更接近真实建模过程。

A001 展示了机理题中“公式后留一句人话”的用法。作者在 RK4 推进式后逐项说明起点、
中点和终点斜率，再用“易知，当上述斜率加权平均时，时间中点斜率权值最大”收束。
先观察前期位移不稳定，才截取 $100$--$180\,\mathrm{s}$ 计算平均功率；阻尼优化先用
增大步长粗略遍历，再用遗传算法求精。旋转阻尼影响小时，只作“可能是纵摇角速度较小”
的局部解释；周期误差和 $\pm1\%$ 转动惯量扰动分别承担检验与可靠性任务。这种
“现象—时间窗口—求解分工—检验规则”比直接写“使用 RK4 和 GA”更接近人类作者推进。

A022 则把“边界先行”写得很清楚：先计算平衡位高 $h_0$ 排除浸没和锥面暴露，再进入运动方程。第二问不是把复数解和数值积分并列，而是由前一问的数值结果授权改用：恒定阻尼可写出解，幂律阻尼则用 \texttt{trapz} 配合粗搜、局部加密和排序。问题三在相对垂荡很小时才把振子放入以浮子为非惯性系的纵摇算式，问题四再由图形发现旋转阻尼单调增长后设为上界并只局部加密纵摇阻尼。“观察—授权—改用—近似—边界”是比“列出 ode45、trapz、扫描”更像人类论文的连接。

A171 把机理题的“人话”写在方法切换之前：先算平衡压缩量和吃水，说明分段浮力在哪些区间适用；常阻尼是常系数方程，故列出拉普拉斯变换的三步，幂律阻尼解析解难求，才改用四阶 Runge--Kutta。优化段先定义瞬时功率、最后十个稳定周期和约束，再给阻抗匹配或梯形积分的结果。每种情形都用 $10,20,40,60,100\,\mathrm{s}$ 交付同一组状态量，读者可以逐项对照；检验段又把初始过渡误差和稳态误差分开，最后落到“可以接受”的判断。可迁移的不是“任何题都用 RK4”，而是“边界—方程性质—方法—同口径结果—判定”的段落顺序。

A0127 把工程计算中最容易被 AI 写空的“为什么这样算”保留下来。问题分析先写“它实际上是一道优化问题”，再把第二问称为第一问的反问题、第三问写成放宽约束；读完题面公式后，又逐项判断哪些量可直接获得，最后才说计算关键是五类光学效率。塔影没有用“为简化计算”一笔带过，而是先交代每面镜姿态不同带来的边界计算负担，再以完全覆盖镜数量占优、镜面中心对称和边界附近统计补偿，推出镜面中心是否落入阴影区的二值规则。后面的空间云图也不是装饰：北侧效率较高，于是镜面向北集中、塔向南移动，塔位变量随之限定在负 $Y$ 轴。搜索段同样先写直接高精度遍历的计算量，再让二分、大步长和小步长分别承担定位区间与提高精度。可迁移的是“问题身份—剩余未知—简化依据—可执行判定—空间证据—计算预算”的自然推进，不是机械复用“考虑到”“因此”等词。

A0165 又展示了另一种更接近人工建模记录的推进。摘要不是先报算法名，而是把定日镜计算拆成太阳位置、镜面法向、遮挡、截断和功率五个传播阶段，分别用“确定”“计算”“判断”和“从而得到”交代接口。几何准备段先说“可将矩形分为两个三角形”，再把点在矩形内的问题等价成两次三角形判定；遮挡段只检查六个近邻，并用 $2\times2$ 栅格构造快速计算核，且报告其与完整计算的偏差小于 $2\%$。优化段分析一个变量时明确固定其余变量，由单峰曲线才采用三分搜索；空间散点图显示北侧效率占优后，塔位不是凭经验设定，而是向南移动。第三问沿用第二问时又把“保留什么”和“新求什么”分开：镜宽、塔位和布场被冻结，只搜索各圈高度。检验也不是再报一个误差数，而是先预言峰值应随太阳视运动移动，主动改换五个采样时刻，再比较散点峰值的迁移方向。值得迁移的是“过程分段—等价转化—快速核及偏差—固定变量—图形触发设计—继承与冻结—结构预言检验”的信息顺序。

A0175 更像一份把程序控制流翻译成论文的人工记录。塔影段先写“若塔身对某入射光线造成阴影遮挡损失，那么不再考虑……”；只有塔影未发生，才继续判断入射阴影、反射遮挡和集热器接收。邻接矩阵也不是突然出现：作者先把 1745 面镜、镜面网格和锥束角度共同造成的计算量写出来，再用遮挡主要发生在邻近镜之间这一物理理由，把候选限制在 20 m 内。步长由 6 m 逐步降到 0.25 m，单位面积输出仍在上升，正文只说从 1.5 m 起“增长的速度变得缓慢”，因而“认为……接近稳定”；这种有限语气比“充分收敛”更符合证据。问题二先以 $10\%$ 镜面重复 10 次估计粗参数，遍历旋转角和高度发现影响不大后冻结，再把预算转向镜宽和塔位，最后用 $30\%$ 镜面重复 5 次汇总。表 7 得到 $68.2443\,\mathrm{MW}$ 后，作者没有停在报数，而是从“额定功率只需要达到即可”继续推理：删去高遮挡镜面会同时降低总面积、总功率和相互遮挡，因此应进一步兼顾单位面积输出。值得迁移的是“条件—停止—继续”“计算量—局部性—邻接”“序列—趋缓—工作步长”“弱影响—冻结—转移预算”和“达标结果—连锁作用—目标修正”。

A092 的自然感来自作者愿意把“还剩什么没算出来”写清楚。把题面数据代入第一问模型后，大气透射、余弦与遮挡等量可以直接获得，正文才说“唯一无法直接求得的是截断效率”，并用随机入射方向、随机镜面点和 10,000 次接收判定解释蒙特卡洛的职责。入射遮挡与反射遮挡可能重叠时，作者也没有直接相加面积，而是先指出重复计数风险，再离散镜面、统计两区域并集内的点。结果图的推进同样分两步：环序号曲线只负责发现北侧高值并写“由此推测”，随后专门选春分日上午 10 时画二维分布，才把空间猜想转成布局证据。第二问固定其余参数，分别沿东西和南北扰动塔位；东西响应关于原点近似对称，上午和下午的影响又相互抵消，因此横向坐标被主动冻结。第三问则从前问的同心圆结果继续推理：同心圆各向同性，同一圈可共享镜面尺寸和高度，逐镜变量便压成逐圈变量，求解方法也由高维 GA 切换为变步长遍历。这里最值得写入 Skill 的不是算法名称，而是“逐项排除—唯一剩余量—随机估计”“曲线—空间猜想—代表截面复核”“方向扰动—对称抵消—冻结变量”和“前问结构—对称依据—按组降维”的段落动作。论文将方案与 100 组随机可行解比较时，结论只应保留在该样本背景内。

A016 的人类感来自作者把数学多解、计算缩域和数值回退都写成了可见判断。相邻把手定长方程有多根时，正文先说板凳依次排列，故角度必须满足前后顺序；又因相邻把手只转过一个小角度，才在可行根中取最近根。碰撞段也不是泛称“降低计算复杂度”，而是先用下游板凳沿前件轨迹运动的事实，判断最早碰撞主要涉及龙头或第二板凳，再列出真正要检查的顶点；可能晚些出现的特殊碰撞，则由“按照时间遍历，找到开始碰撞的时刻立即停止”排除。

逐步逼近和首次碰撞搜索都保留“试一步—越界—退一步—缩步—再判断”，最小步长随后被译成误差小于其两倍。面对题目要求的“调头曲线还能否缩短”，作者没有为了交付一个优化数值而强行搜索，而是利用相切和比例关系证明两段圆弧总长度固定；证明没有自由度后，才继续计算题设比例下的调头路径。问题五同样先证明任一把手速度变化 $k$ 倍会沿链条传递，再整体缩放问题四的速度矩阵。这里值得迁移的是“物理顺序—选根”“最早事件—缩域”“越界—回退—误差”“固定量—不可优化”和“局部齐次—整体缩放”，而不是单独模仿“因此、可以发现、可见”。

\subsubsection{B 类：变量、约束与策略必须逐项对应}

B 类先写决策对象、目标口径和约束来源。典型信息链是“令……表示是否选择……”
“由产能/路径/时序要求得到约束……”“为比较冲突目标，构造……并保留 Pareto
方案”“在相同函数评估预算下，改进方法相对基线……”。每个二元变量应能翻译回
一个实际动作，每个约束都能回到题意，每个输出都能转化为方案表。

不能只给元启发式流程图和最优值。论文需报告编码、初始化、约束修复、参数、
终止条件、随机种子或多次运行统计；能使用精确规划时，还应报告界或间隙。图表
以决策结构、算法流程、收敛/Pareto 图、方案对照和敏感性图为主。

B078、B108、B125、B175、B007、B026、B050、B160、B030、B035、B086、B226、B311、
B477、B159、B195、B196、B060 与 B157 共十九篇 B 题均已通过人工门禁。跨论文反复出现的
写作骨架是“信息集—状态—动作—转移—目标—求解—检验”：
已知信息的单人层先给最短路、动态规划或整数规划基线；未知信息层说明概率来源、
情景生成器、缺货风险和滚动再决策，并把知道完整未来的离线全知上界与只使用当时
信息的在线非预见策略分开；多人层先判定共同收益、一般和、零和或拥塞结构，再定义
联合状态和支付，依次检查最佳反应、支配关系、纯均衡、内点混合解与边界解；若是
合作多目标，则明确加权、词典序、约束法或 Pareto 选择。终点
残值必须按题意进入目标，整数库存和路线候选须覆盖完整可行域；候选集上的混合均衡
不能直接写成完整路径空间的确定最优路线，逐日局部均衡也不能替代跨日状态耦合的全局策略。

相应结果段应把三类证据分开：确定性层报告目标值、库存余量和约束满足；随机层
报告独立重复数、种子、均值、标准差、分位数与失败率，并使类别计数、筛选后分母
和比例在舍入容差内互相还原；多人层报告支付如何映射到题面收益、候选集覆盖和各
玩家可行性。灵敏度分析还须声明每个扰动点是固定原策略
复算还是重新优化，并同时记录目标变化、策略切换和活跃约束。

B007 给出了另一种常见的 B 类数据--优化混合链：先按受控条件整理实验数据，利用
Pearson 相关与二项式拟合辨认单因素关系，再用 Q 型聚类归并催化剂组合；随后在同一
数据和同一指标下比较多元回归与偏最小二乘，针对原模型未直接描述产物收率这一
缺口建立收率回归，并继续加入温度交互项，最后把回归关系送入 LINGO 优化，以均匀
设计补充有限预算下的实验点。这里的优秀之处不是模型多，而是每次增加复杂度前都
先指出上一环节没有回答什么，下一环节又有明确输入和输出。

其结果段也采用容易核读的成对结构：先报告无温度限制下的最优值，再写温度不高于
$350\,^{\circ}\mathrm{C}$ 时的受限最优值，并说明为什么工程上需要第二个方案。实验设计段
先写“现有两种方法可以选择”，逐一说明适用性和舍弃理由，再落到均匀设计及安全
修正。读者看到的是一次真实取舍，而不是事后把唯一方案包装成必然选择。

B026 把这条链继续向“证据缺口驱动实验”推进。作者先用三次样条统一温度水平，
再从散点趋势和转化率上界说明 Logistic 的平台形态；选择性和时间曲线则分别保留
二次关系与分段下降后的稳定段。双因素方差分析先回答温度、催化剂组合是否有影响，
多重比较再定位具体水平，随后把催化剂配方改写为连续变量，分别建立两种装料方式的
二次回归并在全温区、温度不高于 $350\,^{\circ}\mathrm{C}$ 两种口径下求解。这里的
人类写法在于每一步都给出“前一张表还没有回答什么”：缺失组合补测、未观测区间补测、
保持单一变量的对照和最终五次新增实验，分别对应不同的裁决问题；并非笼统声称“增加
实验以验证模型”。

该篇还留下可直接复用的结果句式：先以 F 值和 $p$ 值判断因素是否有影响，再以均值差
回答哪个水平更优；装料方式 I/II 的拟合优度分别为 $90.16\%/90.78\%$，四种优化
情景保持同一目标口径，结果表中的温度为 $400/349/400/349\,^{\circ}\mathrm{C}$。
这些数字只作为论文报告值的页码证据，写作时仍需说明比较条件和结果用途。

B050 又补出两种很像真实建模过程的推进方式。其一，候选模型并不是按拟合优度机械
排序：四次多项式在五个样本点上得到 $R^2=1$，作者随即指出这只是把有限样本完全
插值，无法据此判断整体规律，因而改用三次多项式。这样一段同时留下候选结果、样本
限制和最终裁决，比“综合考虑后选择三次模型”更容易让读者相信取舍确实发生过。

其二，递进回归每次都保留前一轮指标和尚未解决的问题。基础结构的决定系数/调整后
决定系数为 $0.795433/0.685442$，加入非线性变换后变为 $0.828916/0.696456$，再加入
交互项后达到 $0.934994/0.802786$。这些数字不是并排陈列，而是分别承担“现有模型还
缺什么”和“新增结构是否值得”的判断。随后，作者先写多项式回归送入规划后结果不理想，
才改用 BP 网络；方法切换由具体失败触发，不用“为提高精度”一笔带过。最后五次新增
实验先分为三次峰值探索和两次预测确认，再列实验点及用途，读者因此知道每个点是在
扩展搜索范围还是核对模型读数。

B160 把“控制变量法”和“回归效果较好”这两类常被写空的段落补成了可核读结构。
控制变量比较并不只报因素名称，而是先列其余保持条件和唯一改变量，点名可比组，
随后在同一温度下分别读取乙醇转化率与 C4 选择性，最后只收束当前组的局部结论。
因此，读者既能核对组间是否可比，也不会因单个响应改善就误把配方写成总体占优。

回归段则保留“初始模型--诊断图--多重共线性--逐步回归--修正模型”的真实轨迹。
转化率模型依次报告 $R=0.91$、$R^2=0.832$、$RMSE=0.0417$、$F=39.74$ 与
$p=6.17\times10^{-27}$；选择性模型相应为 $0.89$、$0.80$、$0.0453$、$36.03$ 与
$1.14\times10^{-24}$。作者把前两项称作相对效果，把 RMSE 作为绝对效果，再由 $F/p$
说明整体显著性，三层证据各有任务。问题三又按全温区和不高于 $350\,^{\circ}\mathrm{C}$
重新求解，论文报告收率分别为 $56.3\%$ 与 $29.8\%$；D1--D5 则分别承担曲线延拓、
缺失温区、两类响应水平组合和优化点实测的裁决，不以一句“增加五次实验验证模型”
掩去用途差异。

B030 的写法从“先求全”开始。论文先把三架无人机的相对位置列成三类九种形态，
随后才写“在求解过程中，发现”其中两类可通过旋转重叠，于是把重复推导压成六个
定位分支。这里的简化不是作者预先宣告的愿望，而是完整分类后观察到的几何事实。
算法段同样不写空泛的“进行仿真”：先固定发射机和极径，再列另一发射机、接收机、
六种模型及两个 $[-1^\circ,1^\circ]$ 误差网格，最后写“记录所得全部结果，并绘图输出”。
对象、范围和产物都在同一段内闭合。

该篇还给出两种长文中很自然的节制。10500 组结果不硬塞进正文，而是先说明完整结果
保存于 Excel 支撑材料，正文只展示一个代表情形；锥形编队也不重新铺写整套公式，而是
先以 FY05 为圆心转成半径 $50\,\mathrm{m}$、$50\sqrt{3}\,\mathrm{m}$ 和
$100\,\mathrm{m}$ 的三层圆周，再写“其调整推导与问题一（3）完全一致，对此不再重述”。
结果段先报最大角度和半径调整量，继而用飞行距离解释两者为何量级不同，再以“虽有误差，
但误差较小”保留让步，最后作整体坐标对照。这样的行文比连续使用“有效地”“充分证明”
更像真实作者，因为省略、复用和判断都有可指出的依据。

B035 让另一种方法切换显得更像真实求解。论文先写极坐标正弦定理的四区域分类，
继而直说“这种分类讨论求解方程的做法过于麻烦”，再以定弦定角对应两圆轨迹改写定位。
新方法并非只换一个名称：旋转矩阵与叉积判定圆心所在侧，两圆心连线给出另一交点的
对称关系，原先的分区方程因此变成可画、可检查的几何构造。这里可迁移的是“旧方法--
可见负担--保留不变量--新构造--减少什么计算”的完整转折。

误差论证也没有用一句“经证明可识别”带过。作者先给一个具体偏差例，再把候选位置
限制到轨迹圆与极角扇形相交的短弧；随后检查端点，证明第三个角的变化范围不超过
$2^\circ$，最后用分区角度表穷举无交叠区间。调整算法比较时则承认“在局部上可能
贪心算法占优势”，再解释前视搜索为何会先修正误差更大的点，使全局目标反转。六轮
方案的 $L_2$ 损失约为 $10^{-23}\,\mathrm{m}^2$，但正文随即说明实际应用只取前五轮、
共 16 架次；额外数值精度并不自动值得通信与编队代价。问题二复用前问前又集中列出
三点共线、四点共圆等退化情形及其修复。这些段落把让步、停止和失败接口都留在正文中。

B086 进一步说明，几何优化论文的“信息合法性”也可以写得清楚而自然。论文先由
“无人机间信息不共享”排除直接使用全局队形误差，再让每架接收机只比较自身可见的
角度向量与理想角度向量。随后又定义全局误差，但紧接着说明该量只在仿真结束后用于
评价，各无人机并没有据此调整。于是“局部目标”和“全局评价”不是两个容易混淆的
公式，而是分别回答“现场能用什么信息”和“赛后怎样比较方案”。

径向预处理也没有停在 $O'_j/R\leq1.049$ 这一误差界上。作者先据此把题设
$\pm15\%$ 的半径偏差压到约 $\pm5\%$，再用“更重要的是”说明搜索区域缩小后，局部
穷举需要的信号与计算次数也随之减少。两种调整方案使用同一组原始坐标，先比较收敛
速度，再比较稳定误差，最后估算搜索发射代理：方案一和方案二分别为 $134400$ 与
$25600$，后者约为前者的 $19\%$。论文又特别指出该代理没有实际调整次数的直接含义，
因此最终裁决同时保留了效果、资源和口径边界。

锥形编队段把三条几何引理写成了实际流程。每条引理都交代固定点、调整点、控制角或
比例以及目标图形，随后映射到 A1--A15 的具体机群，组成三阶段调整。平面拟合检验值
由 $3.893$ 降至 $0.206$。稳健性段则先承认单组原始坐标不足以检验模型，再给半径和
角度加入高斯微扰，使用十组数据重复计算，并把结论限定在论文所报 $10^{-7}$--
$10^{-5}$ 收敛量级内。这样的承认、补证和限定，比直接写“模型具有较强鲁棒性”更像
真实作者的论证。

B226 又补出一类很少被算法清单捕捉到的人工写法。作者先说“由问题一知，只要已知
水深和测线与坡面法线投影的夹角，就可以求覆盖宽度”，于是三维问题没有被包装成
一个新的“综合模型”，而是被压成两个待求量。随后用“为了让模型的建立过程更加
直观易懂，我们记……为面 I，……为面 II”给几何对象命名，才进入法向量、叉乘和
线面夹角。这样的段落先让读者知道在算什么，再让代数承担计算，比直接堆向量式自然。

测线方向也没有只靠规范背书。正文先引用作业规程说明主测线通常平行等深线，随即
写“出于建模的严谨性，我们在下文证明”，再用命题和同长度候选的梯形/矩形覆盖面积
比较完成授权。测线间距则从东边界锚定第一条测线，以 1 m 为步长向西搜索，只保留
$10\%$--$20\%$ 重叠率并取最低可行值。这里的“以……为始边”“恰好能达到……”“从
第一条测线开始”不是装饰性连接词，而是直接对应初始化、步长、可行域和选择规则。

第四问保留了另一种真实取舍。作者直说“受到问题三模型的启发，我们人为地将海域
划分为四个区域”，并紧跟等深线方向这一依据；局部平面拟合后，区域 III 和 IV 的
均方误差过大，才只拆这两个失败区域，已合格区域保持不动。结果段也不写统一的
“方案合理”：区域 V 因坡度大、单线覆盖窄而测线多；区域 II 因沿用区域 I 的延长线
以降低作业难度而出现集中漏测。异常区域、局部原因和数值后果逐项相接，正是应写入
Skill 的人类段落组织。

B311 对“题意存在两种理解”给出了一条更值得保留的写作路径。作者没有先替题面
选定定义，而是分别建立两个重叠率模型，把同一组测线数值代入，再问两种模型会在
CAD 图中留下什么不同的可观察区域。图形裁决后才舍弃模型一、保留模型二。这样的
“歧义--备选--共同例子--判别量--取舍”把真实犹豫留在正文中，比直接写“经分析采用
定义二”更像人类的求解记录。

跨问推广也只改真正变化的量。作者先比较本问与问题一的几何结构，保留原覆盖关系，
再重新计算当前位置水深和测线方向与坡面的夹角；函数接口因此只替换这两个输入。
方向判断没有因为坡度小便立刻忽略，而是把它积累到 2 海里尺度，发现两端覆盖宽度
已经相差数百米、可行间距区间不能兼容后，才改用平行等深线。`GB12327-2022` 出现在
定量裁决之后，只作补充印证。这种句序先交代题内证据，再引外部规范，不让标准替代推导。

模型检查与递推收尾同样具体。“首先、其次、再次、最后”分别核对对称性、单调性、
沿等深线方向的不变性和特殊角度下的极限式，并非为了凑四句排比。问题三从西边界
确定首线，按 $10\%$ 重叠递推到第 34 条，再用东侧剩余 $22.1\,\mathrm{m}$ 和
$3713.7>3704$ 说明第 35 条应舍弃；问题四又以等深线间距最大的最平缓位置作为最不利点，
由该处完全覆盖推出其余位置覆盖。结构预言、边界余量和最不利位置各自承担一个结论，
所以段落虽使用常见连接词，仍有清楚的人类判断痕迹。

B477 把这种真实判断继续推进到“定义失效--换法--局部返工”三个节点。旧重叠率并非
被一句“存在局限”带过：作者固定相同测线间距和重叠长度，画出直观重叠更大的情形
反而得到更小比例，写明继续使用会使指标失去意义，随后只把间距改成坡面方向上的
投影量。二维关系推广到三维以前，又先证明航向固定时局部坡度唯一，再把当前位置
深度和等效坡度送回原式。段落的自然感来自反例只触发必要修改，证明也确实授权下一步。

真实海域部分没有用“为提高精度，采用随机森林”替代选择过程。论文先把纯离散加密
的负担拆成精度、局部失效、误差累积以及数据量与耗时四项，再要求学习模块输出任意点
的深度、梯度、坡度和覆盖宽度，因而读者知道新方法交付给哪一段路线算法。第一次搜索
出现起点依赖、局部漏测和可用性差时，正文也保留了这个失败状态；后续按等深线分七区，
只对浅而平缓的 I、II 区另行拟合。这里应模仿的不是“针对上述问题”六个字，而是旧方案
的具体负担、新方案的下游接口以及失败区域和修复范围都没有省略。

结果交付同样按阅读任务分流。路线过密时，正文图只说明位置，精确坐标移到附件，
主文继续按区给条数、长度与面积，并汇总 94 条、$434662\,\mathrm{m}$、有效覆盖
$99.956\%$ 和漏测 $0.04\%$。多起始线只支持已比较候选中的代表；连续航线和返航仍留给
Hamilton/Euler 路线接口。写“取最优方案”以前先交代候选范围，正是人类论文避免事后
夸大结论的一种朴素写法。

B159 与 B195 把抽检题里的选择过程写得很有层次。B159 先问企业更怕把次品放行，还是
更怕把合格批拒收，由这一态度反向确定 $H_0/H_1$，再从功效、样本量和成本进入决策；
B195 面对没有真实抽检样本的条件，先公开模拟器及其假设，待每一步费用、时间和状态转移
闭合后，才压缩重复流程。两篇都不把“采用动态规划”当作起点，而让检验立场、数据条件和
流程复杂度逐级授权后面的模型。

B196 的段落从结构计数起步：节点数、反馈边数和候选动作一旦达到某个规模，穷举负担才
成为换用蚁群或遗传算法的理由。结果解释也没有停在收敛曲线，而是先找反复出现的关键节点，
再说明这些节点对应哪一步真实操作；扰动后则依次回答哪些节点发生变化、总成本怎样变化、
为什么是这一方向。这样的“结构--搜索--热点--执行”比列出三个智能优化算法更有内容。

B060 与 B157 代表另一条机理反演写法。B060 把局部谱线与全曲线反演分工：前者便于锁定
候选区间，后者承担最终参数估计；若中间频谱发生变化，还要把差异重新送入厚度模型，不能
停在中间变量。B157 则先说明频率干扰会怎样影响厚度反演，再决定分段修正；当主体厚度
几乎不变而局部残差下降时，作者把贡献限定在局部拟合，不把有限改善扩大成“全面提升”。
这种主动收窄结论强度的语气，是人类技术论文很重要的一种可信感。

\subsubsection{C 类：数据口径、隔离验证与决策解释并重}

C 类从字段、样本单位、时间范围、缺失异常和训练/验证边界写起。有效表达包括
“以……为一条观测，按主体/时间划分训练与验证集”“所有标准化和特征选择仅在
训练折拟合”“相对持续性/线性基线，模型在……指标上……”“该变化主要对应……
特征，但结论限于……范围”。评价排序需解释指标正向化、权重来源和排序稳定性；
预测需把误差转换为补货、时点或风险决策的影响。

需要避开的是全数据清洗后再切分、用随机切分检验未来预测、把相关性写成因果、
只报告最优模型以及用二维降维图证明分类性能。图表以数据质量、分布与相关、
验证协议、残差/校准、特征影响和决策结果为主。

C006、C085、C169、C283、C155、C229、C050、C126、C228、C109、C142、C170、C227、
C305、C038、C063、C094、C234、C023 与 C132 共二十篇 C 题均已通过人工门禁。以下归纳
同时覆盖业务统计、规划决策和医学筛查，但每种句式仍须回到本题的样本单位、时间口径和
决策对象，不能把不同应用中的模型名机械拼接。
C006 的正向价值首先体现在“把业务语言改写成可计算量”的过程。论文没有直接把“供货
连续性”当成一个无来由的分数，而是拆成间隔次数、平均间隔周数和平均连续供货周数，
再说明各指标的成本/效益方向、归一化和熵权。这里可迁移的不是熵权--TOPSIS 这个名称，
而是“业务概念--不可直接观测--代理指标--方向--权重--解释”的完整段落。

该篇在周供货能力处还主动拒绝了更省事的历史均值：均值会把不同周次抹平，经济作物的
产量又具有周期性，故改为在十个 24 周周期中取相同周次的历史位置，再用 GP 估计订货量
对应的供货误差。逐周订购时，上周库存、实际供货和消耗量共同形成下一周初始状态；
二元供应商筛选与逐周订购使用遗传算法，而转运环节按损耗率、整批供应、剩余运力和 A 类
优先等规则分配。这样的写法把三个求解角色分开，避免只因同属一条业务链便统称“用 GA
优化”。

结果解释也保留了人的判断痕迹。作者先给 24 周总体方案，再选第 3 周和第 13 周说明库存、
供货与转运怎样在具体周次闭合；发现原方案没有表现出材料偏好后，先诊断成本系数所对应的
替代机制，再只调整 A/C 原料的等产能约束。调整后的成本最大约降 $5\%$、平均仅降
$0.956\%$，论文明确承认改善不明显，并把原因落到供应商数量少、供货能力贴近需求和可行域
狭窄，而不是把正向变化包装成“显著提升”。问题四沿用转运模型时，又明确说明输入、输出和
优化目标均未改变；这比一句“沿用上问模型”更能说明模块为何可复用。

C085 把另一类业务判断写得更具体。四个供应商指标不是只给名称和权重，而是逐项解释
“统计量怎样定义、数值方向表示什么、它对应哪一种供货能力”；稳定供货量又先受两条
规则筛选：供货/订货超过 $0.9$，且供货量不超过同周期位置均值的三倍，只有同时满足时
才取最大候选值。这样的段落先让阈值承担稳定性与异常排除作用，再给期望量，不以
“经预处理得到”省略判断过程。

算法改进也由可指认的失败触发。固定周权重可能使总体供货看似充足，却让个别周严重
缺货，所以作者按当前供货与期望产能的缺口更新下一轮权重；第 7、8 周在数据中连续处于
淡季，因而前视窗口取两周并提前备料。转运部分则先拆分超过单个转运商容量的大项，再按
可靠性和材料价值做 0--1 背包，最后对未装入项贪心补救。结果比较时，历史采购成本虽低，
但历史平均产能没有达到 28200，故改用采购成本/产能比较 $0.722032258$ 与 $0.721501175$。
这几处都先写“原口径或原算法具体在哪儿失真”，再给修正动作；比“综合考虑多种因素，
采用混合算法优化方案”更接近人类作者的真实推理。

C169 展示了另一种很容易被 AI 文本抹掉的动作：作者会说明哪些观察\emph{不进入}模型。
前 50 家供应商的序列中既有周期性异常，也有非周期尖峰；论文先解释两类现象的可能成因，
随后用“尖峰会干扰周期规律”和“新订货方案会打破旧规律”两个理由，明确下文不单独建模。
类似地，初始库存没有方便地设为 0，而是以“企业并非刚刚成立”为起点，回到历史 24 周
周期的同位状态估计。转运模型也先判断损耗无记忆，确认过去损耗不影响本周后，才按周拆分；
这种次序比“为简化计算，将问题分解为 24 个子问题”更可信。

该篇还把业务后果直接塑造成数学结构。缺供对生产的危害重于超供，作者先列分段函数应满足
的符号、曲率和两侧大小条件，再构造 $e^x-1$ 与 $e^{-2x}-1$，而不是先选函数再补解释。
候选集的前 50 家又由“覆盖总供货量 90\% 以上、尾部仅个位数、继续增加会损害经济性”共同
限定。跨周聚合时，最少供应商取 24 个逐周最小数的最大值，最大产能取 24 个逐周上界的
最小值；聚合方向不同，是因为前者要找统一可用的供应商数，后者要保证每周都能达到产能。

C283 把评价指标进一步写成附件上的实际操作。忙期均量并非对 240 周直接求平均，而是先去掉
无订单形成的零周；信誉度又让相对订供偏差乘以订货量的对数，使大额且按量履约的订单提供更多
证据。熵权结果出现供给弹性等业务重要指标权重过低的情形后，作者没有用一句“主客观结合”
带过，而是先把指标归到绩效、战略意义和潜力，再在适度范围内用 AHP 修正并通过一致性门。
这形成“数据权重--具体异常--业务冲突--有限修正--一致性检查”的完整段落。

双层贪心的选型和停止也保留了真实取舍。论文先写线性规划的变量/约束关系，再指出动态规划的
状态空间约为 $6000\times50\times24$，随后利用供应商和转运商已有排序构造两层贪心。库存达到
两周需求后立即停止，不是因为程序“收敛”，而是继续下单会转向较差供应方或运输方并增加存运
成本。四个关键订货式共用历史忙期均量，再按约束、成本、初期稳健性和现实含义逐个淘汰；随机
运行则明确只取“频率较高且目标较优”的典型代表，不把一次输出包装成唯一最优。

问题三的类别权重同时写了上下边界：过小不能突出 A/C 的存储差异，过大又会牺牲采购经济性和
合作可靠性。到问题四，目标由成本/损耗最小转为产能最大，先前排名靠后的供应商也可能贡献供货，
所以候选从前 50 家扩到全部 176 家。结果段发现实际停止条件与预期瓶颈不完全一致时，作者保留
疑问并给出“提高进货参数”或“降低偏差率”两种解释；生产曲线与接收曲线同向、约滞后 3 周，
则把库存传导写成了可读的时间关系。

C155 先处理“缺失到底表示什么”，而不是把所有空格交给同一种插补。化学成分空白按题意表示
未检出，故填为 0；颜色字段却不能照此处理，因为相同已知条件仍可能对应多种颜色，且表面随机
取样不足以代表整件文物。论文据此明确拒绝猜测性填补，只在涉及颜色的分析中忽略 19、40、48、
58 号样品。类似地，成分总和低于 $85\%$ 的 15、17 号点先由题面阈值判为无效，再剔除；其余
样本闭合到 $100\%$ 后才作中心化对数比变换。这种“题面语义--具体编号--判断--动作--影响范围”
比一句“对缺失值和异常值进行预处理”更能让读者知道数据发生了什么。

统计方法也由前提而不是方法名分流。作者先列联表和期望计数，满足 Pearson 卡方检验前提的组
直接使用 Pearson，不满足的组改用 Yates 校正；风化与类型的 $p=0.009$，与颜色、纹饰的
$p$ 分别为 $0.507$、$0.084$。随后 $k=2$ 的 Q 型聚类借助前问趋势和已知标签命名风化/未风化
簇，又用聚类结果反向核对前问判断；$k=4$ 再把风化展开为未风化、轻度、中度和重度四阶段。
同一条证据在前后问题之间双向使用，使各问不是互不相干的算法堆叠。

预测段先逐成分比较二次回归的 $R^2$，只预测拟合较好的成分，对较低者采取不变策略；若异常点
使“曲线到观测值的距离”被放大，则只对这个唯一敏感量作单调压缩，再经逆对数比和闭合约束恢复
总和为 $100\%$ 的成分向量。分类段又把 R 型和 Q 型聚类的职责拆开：R 型先将十四个相关成分
变量分组并选代表量，Q 型再用代表量划分样品亚类。未知样品由两棵按风化状态分层的 PbO 决策树
分类，并以 Q 聚类交叉核对；最后固定每一种成分为灰色关联母序列，再比较两类玻璃的最大、最小
关联及分布集中程度。这里可迁移的是“指标门控--局部修正--闭合验证”“变量先降相关、样本后分组”
以及“同一基准下比较极值和分布”，不是把所有方法名称并列到摘要中。

C229 进一步说明，即使缺失的是同一个颜色字段，填补方法也不必一刀切。论文先用纹饰、表面
风化与玻璃类型组成匹配条件，再逐个核对候选样品的参考质量。40、58 号分别能找到成分相近且
可参照的 39、11 号，因此用热卡填为深绿、浅蓝；19、48 号所在组虽有匹配对象，但多数化学
样本来自未风化采样点，参考性较低，故转用众数并填为浅蓝。正文把“分组条件--候选编号--
参考质量--方法分支--填补结果”连续写出，比笼统声称“对定性变量采用合理插补”更可信。

化学成分的空白与数值 0 也先按数据生成过程解释：前者来自低于检测限，后者可能是保留位数后的
舍入结果，二者都不能当作普通零值进入对数比。作者以列最小正值近似检测限，取
$\delta_j=0.65e_j$ 作乘法替换，再把向量闭合到 $100\%$。SiO2 占比过大而遮蔽其他成分时，
论文保留完整成分图，又补一张去除 SiO2 的图，最后写“结合两图进行分析”。这种做法没有为了
看清小量而删掉主量证据，也没有让一张比例失衡的图承担全部结论。

分类段先用决策树给出可直接复核的 PbO=1.5 阈值，在 47 个训练样品和 20 个测试样品上报告
测试准确率 $100\%$；随后 PLS-DA 以 PbO、K2O、BaO、SrO 的 VIP 值复核多变量差异。亚类
没有已知标签，作者据此明确把任务判为无监督学习，再由 WSS 曲线选 $k=3$。铅钡类三簇较清楚，
高钾类没有明显分界，正文只写“仅可粗略将其分为三类”。这里可迁移的是“可解释规则--多变量
复核”和“无标签身份--诚实限定聚类强度”，而不是把两种模型或两类图统一包装成效果良好。

未知样品的 PLS 维数同时比较十折随机段交叉验证误差与累计解释率：三个成分时 CV RMSEP 为
$0.1371$，继续增加到四个成分只到 $0.1372$，而 X 与类别累计解释率已达 $93.55\%$ 与
$92.00\%$，故停止在三维。A1、A6、A7 判为高钾，其余五件判为铅钡，树与 PLS 结论一致。
问题四又从数据结构推出检验：相关矩阵对称，故只取上三角；相同成分对一一对应，故用配对
Wilcoxon，$p=0.905$。参数选择和检验选择都由可检查的两层理由推进，而非由方法名称代替判断。

C050 把数据侧写写成模型入口，而不是把探索性图表集中摆放后便不再使用。论文先从品种丰富度、
季度销量、损耗、描述统计和时序变化观察附件，并在负价记录的业务含义明确后剔除退货数据；后续
相关分析、半年预测窗口、可售单品范围和优化约束都能回指这些观察。单品数量过多时，作者也不是
随手挑几条曲线，而是先规定各类展示总销量最高者，以及总销量最低但不少于 $2000\,\mathrm{kg}$
者，其余结果转入附件。这样的顺序形成“侧写对象--实际观察--下游判断”和“代表规则--正文展示--
全量附件”的双重接口，使图表选择本身可以复核。

回归和人工特征的解释继续落回业务口径。论文没有停在相关系数、显著性或线性方程，而是逐类说明
销量每增加 $1\,\mathrm{kg}$ 时成本加成定价相应增加或减少多少元。问题四又分别以品类单品种数、
季度中心附近的节气和主要节日构造丰富度、季节与节日指标，并先说明代表对象、计算口径及其所代理
的经营现象，再进入灰色关联。可迁移的是“系数--业务单位--方向--决策含义”和“对象--口径--代理
语义--下游用途”，不是在结果段重复一遍公式或把所有人工指标笼统称为特征工程。

该篇还把现有附件能够支持的结论与尚待采集的信息分开。附件内的销量、价格和日期先完成三项关联
分析；客流、竞品、反馈、天气、供应链等信息则另设段落，按“如何采集--能够区分什么机制--支持
何种经营动作”展开。打折频率较低既可能表示商品受欢迎或供不应求，也可能只是促销必要性较弱；
频率较高则可能来自销售乏力或保鲜期较短，故需联查销量、反馈和竞争情况后再决定定价或促销。
这里的关键不是列更多外部变量，而是保留“代理量--竞争机制--追加证据--分情形动作”的推理过程。
正文只保留代表图表、关键数值、选择规则和解释，完整单品表、补充图与代码进入附录，使证据分流与
版面分工保持一致。

C126 的第一条写作价值，是把方法切换前的压力写全。日销售量不服从正态、零值多、三年序列有
季节结构，OpenBUGS 的 MCMC 又有实际计算负担；作者因此不是直接宣布“采用 Bayesian 方法”，
而是先用同期平均把三年数据压成一年，再交代新口径怎样进入先验、迭代、燃烧期和后验诊断。
可迁移的结构是“原数据现象--旧口径的具体困难--转换规则--新输入--求解设置”，它允许读者看见
为什么换模型，而不是用“提高精度、降低复杂度”两个空泛评价代替选择过程。

第二条价值来自业务事实对模型结构的授权。论文先写蔬菜当日未售后隔日大多难以原价销售，由此
把有效记忆限定为两天，再选择二阶 VAR；模型阶数不是孤立的调参结果。预测出现负数时，作者也没有
机械删去，而把它解释为未售且未退的余货、次日可折价以及当日无需补货。这样，“业务事实--时间
尺度--模型阶数”和“数学读数--库存状态--经营动作”各自形成完整自然段，比“综合考虑实际情况”
更能留下人的判断痕迹。

第三条价值是信息压缩和证据分流。相关系数、阈值与显著性先被压成 T/F 判断矩阵；13 维或更大的
结果难以放入正文时，作者明确以花叶类或 7 个水生根茎类为例，正文保留判断规则、代表参数表、
MCMC 诊断和模型评价，完整图表、四组 BUGS 数据块、EViews 工程、R 检验与两问 LINGO 代码进入
支撑材料。VAR 路线筛得 29 个单品后，又从半小时销售频率独立筛选并得到近似名单；这一动作只称
有限交叉核对。预测量随后进入利润收入项，成本、损耗和利润率进入约束，问题三再把最近 7 日反推
进货量、平均计划量、前日余货、新补货和 $2.5\,\mathrm{kg}$ 陈列下限逐项接起。模型之间因此有
可见中间量，而不是几个算法名称排成一条箭头。

C228 把开放任务的第一步写成“先决定看什么”。题目没有指定具体探究方向，作者没有用“从多个
维度深入分析”一笔带过，而是把品类和单品分别放到时间分布、销售分布及相互关系三个观察面上，
并说明这些角度合在一起怎样覆盖销售规律。年、周、日周期的统计读数之后又接上采收季、春节、
周末采购和上下班节奏；图形不作孤立展示，随后用于确定预测窗口，并解释相应的经营现象。

同篇还保留了两次不同性质的方法退让。第一次是试过 ARIMA、回归和灰色预测后，明确写出当前
序列上的效果不理想，再改用 LSTM；第二次是面对近七日单品的缺失、短时无明显趋势和随机波动，
反而不再强行使用复杂预测，而以加权平均给出当前时窗内的估计。前者是“候选--表现--换模”，
后者是“数据能力--复杂方法不再合适--主动降级--限定时窗”，两种取舍都比事后罗列算法优点更接近
真实建模过程。

预测到优化的接口也没有被一句“代入模型”吞掉。LSTM 先给基线销量，双对数弹性再按价格变化
修正需求，修正量进入收入项，成本、损耗和补货量分别进入目标或约束；问题三新增是否进货的
$0/1$ 变量、连续加成率和补货量以及双目标，作者因变量结构改变才从连续 PSO 转向混合编码的
NSGA-II。结果段又把方案放回同季历史日利润区间比较。最后的补采建议从成本导向定价的边界长出，
具体指向竞品价格、消费者预期和收入，而不是另列一份与前文无关的“更多影响因素”。

C109、C142、C170、C227 与 C305 都把企业画像、风险估计和信贷方案串在一条决策链上。
这组论文最值得保留的不是某个分类器名称，而是分问之间的接口意识：发票先整理成企业级特征，
特征产生风险或等级，风险再进入准入、额度、利率和冲击重算。写作时应逐项说明“上一问的哪个
输出成为本问的哪个输入”，并让长表和附录代码承担逐企业追踪，正文只保留规则、代表结果和解释。

这五篇的自然行文还体现在业务量会被重新翻译。交易连续性不是抽象分数，而要落到交易次数、
金额、波动或中断；外部冲击不是一句“考虑疫情影响”，而要落到行业、地区、需求或偿付能力的
参数变化。段落宜写成“由附件可计算……，故以……刻画……；该量进入……项，从而影响……”。
如果只能给单一情景，就把结论限定在该情景，不用“具有普遍适用性”之类空泛收尾。

C038、C063、C094 与 C234 展示了种植规划中四种互补的表达。C038 把七年规划先写成明确的
地块、季次、作物和轮作约束，再用 CVaR 讨论尾部风险；C063 让典型情景负责展示候选差异，
却把最终裁决放到共同随机环境中；C094 将“优先利用、稳定种植、兼顾收益”翻译为管理参数
$p/q$；C234 先从数据表确定索引，再把三类真实邻接关系写进染色体和修复算子。共同环境复评、
管理语言参数化和可行编码，都是“现实判断怎样进入模型”的具体答案。

这组规划论文的结果段也不把“收益更高”单独拿出来。作者会同时写期望收益、最坏情形或尾部
风险，说明为了稳健性牺牲了多少利润；候选差别很小时，再回到活跃约束、种植集中度和实际执行
难度裁决。适合复用的句序是：“方案……在基准情景下……；置于相同的……情景后，……；
差异主要来自……约束，因此选择……。”三个省略处都必须填入本题对象和读数。

C023 把医学筛查中的统计模型与行动后果分开。文章先用混合效应模型处理重复测量，再把达标概率
送入分组和检测时点决策；安全、及时与可执行分别对应不同损失，概率拟合好并不自动等于决策好。
“结果可疑”也被翻译为复测或升级处置。这里可直接迁移的组织是“概率估计--阈值含义--错误代价
--行动规则--稳健核对”，而不是把 Logit、动态规划和 FDR 并排罗列。

C132 先把多次检测记录从“行”还原为“孕妇”，再定义最早达标时点，这是其全文最关键的口径
转换。缺测时按实测证据、保守估计和无法判断三级回退；首次就诊即可获得的静态特征交给 MLP，
随孕周变化的序列交给 LSTM，存在右删失的时点风险交给 DeepHit。作者没有用“多模型融合”
概括这些角色，而是逐一说明各模型接住哪类信息和哪种缺失。

该篇对结果突变的解释也值得保留：当达标概率接近 $90\%$ 阈值时，继续等待的边际代价突然增大，
于是曲线变陡；最终融合不是把各专家在全数据上的输出直接重用，而以折外概率作为元分类器输入。
结果段同时报告总体表现、关键少数类召回和对应样本数，使“整体不错”不会遮住真正需要识别的
高风险人群。
\subsection{十四类章节的写法}

\begin{longtable}{p{1.8cm}p{4.2cm}p{4.0cm}p{4.0cm}}
\caption{章节功能、段落套路与质量门}\label{tab:sections}\\
\toprule
章节 & 必答问题与段落结构 & 推荐表达及图表 & 质量门与常见错误\\
\midrule
\endfirsthead
\toprule
章节 & 必答问题与段落结构 & 推荐表达及图表 & 质量门与常见错误\\
\midrule
\endhead
摘要 & 逐问写对象、模型、关键结果和检验；最后给关键词。 & “针对……，建立……；求得……；经……检验……”；通常不放图。 & 每问均出现，数字与正文一致；不堆算法名，不写“效果良好”。\\
问题重述 & 重组给定量、输出、硬约束、附件和问间依赖。 & “给定……，要求在……约束下确定……”；复杂几何可重绘示意图。 & 不大段照抄，不漏单位、时段或提交文件，不提前引入模型。\\
问题分析 & 每问说明核心矛盾、输入输出、基线、改进和数据流。 & “由于……，故先……，再……”；节末可放总流程图。 & 必须解释选择理由和近邻方法边界，不能只列算法。\\
模型假设 & 一条一义写对象、范围、简化理由和方程后果。 & “在……尺度内忽略……，因此……”；必要时放情景表。 & 题面事实不是假设；不得写“数据真实可靠”；后文可追溯。\\
符号说明 & 定义状态量、决策量、参数、集合、索引、单位和域。 & 三线表分组；局部临时量在首次出现处定义。 & 同符号不跨问改义；公式、代码和附件字段一致。\\
模型建立 & 依次给依据、变量、数学表达、约束、初边值和参数。 & “令……；由……可得……”；结构图靠近定义。 & 公式前有来源、后有解释；约束方向、量纲和变量域正确。\\
模型求解 & 写预处理、初始化、求解器/迭代、参数、停止和输出。 & 按执行顺序写；伪代码和流程图不替代参数表。 & 第三方可复算；不以“调用 MATLAB/Python”结束。\\
结果分析 & 读关键值，比较基线，解释原因，说明约束和决策含义。 & “在……条件下，……由……变为……，表明……”；图表后立即解释。 & 不逐格念表；有单位、有效数字和不确定范围；每问均回应。\\
模型检验 & 明确检验对象、方法、阈值、独立数据、指标和裁决。 & 回代、守恒、残差、网格收敛、交叉验证或消融。 & 不能用训练拟合代替外部能力；失败项也要报告。\\
灵敏度分析 & 定义扰动参数、范围、采样、重复求解和响应指标。 & “当……位于……时结论不变；超过……后策略切换”；放曲线/情景表。 & 扰动后必须重算；区分数值稳定和决策稳定。\\
模型评价 & 将优势绑定机制，将不足绑定假设、数据或算法。 & “由于显式保留……，模型能够……；受……限制……” & 不写任何模型都适用的“简单准确、适用面广”。\\
改进方案 & 对每项局限给新数据/变量/约束、实现方式和验证协议。 & 按“缺陷—改法—验证”对应；可用短表。 & 增加复杂度不是天然改进；方案必须可执行和可检验。\\
参考文献 & 只列正文实际使用的理论、数据、软件与标准来源。 & 统一著录；定义和算法优先原始论文、教材、标准或官方资料。 & 正文引用与条目双向对应；网页含机构、日期和访问日期。\\
附录 & 给入口脚本、环境、参数、随机种子、数据字典和补充结果。 & 文件清单、目录树、可复制代码；主文明确引用。 & 关键结果可由附录复现；不放截图代码，不泄露个人信息。\\
\bottomrule
\end{longtable}

\subsection{可复用表达与“避雷”边界}

可复用的是信息结构，不是整句套用。摘要可用“针对对象和任务—建立核心模型—
给出关键结果—报告检验”的四步；模型段可用“关系依据—变量定义—方程/约束—
式后解释”；结果段可用“读数—比较—机制—决策”；评价段可用“机制带来的能力—
假设造成的边界—可验证改进”。“首先、其次、最后”“通过软件求解可得”“结果
较为理想”“模型具有较强普适性”等句子若没有对象、数值、条件和证据，应删除
或改写。

表达复用不等于隐藏不确定性。OCR 无法确认的字词标为未确认；跨论文统计只报告
检测定义和样本数；DOCX 中公式对象或附件字段若未被结构化提取，则回看原件，不
依据评述补字。逐篇学习首先提取论文为什么读起来像真实研究：作者看见了什么、
为何作出选择、怎样把公式接到结果、又怎样用图表完成判断。只有会改变结论强度或
妨碍复现的边界才进入质量门；个别疏漏不作为主要学习对象，更不以挑错数量衡量
审读深度。

\subsection{人类作者的惯用词、段落骨架与行文节奏}

从 A217、A001、A022、A171、A0127、A0165、A0175、A092、A016、A053、A163、A178、
A242、A196、B007、B026、B050、B160、B030、B035、B086、B226、B311、B477、B159、
B195、B196、B060、B157、C006、C085、C169、C283、C155、C229、C050、C126、C228、
C038、C063、C094、C234、C023 与 C132 共四十四篇逐页细粒度文风样本，并结合其余
十五篇结构化动作卡看，“有人味”不等于
口语化，也不等于把连接词替换一遍。
它来自作者在段落里留下真实的观察和取舍。A217 常用“注意到……，故……”把几何
特征接到降维推导，用“考虑到……，本文将……”说明为何采用某种坐标或计算口径；
B007 常以“观察附件……发现……，利用……判断……”交代数据处理起点，以“由上述
分析可知……，故……”从比较结果转入模型改进。这样的“故”前面必须有可检查的
事实，后面必须落到一个实际动作，不能只连接两句结论。
B026 还常把“但由于……所以我们预测……因此，我们建立……”用于散点和机理上界
之间的推理，并用“但经数据比对，我们发现题目所给出的数据中没有……”引出补测；
这类长承接只在前后事实确实相反或存在缺格时使用。
B050 的常用动作不是增加连接词，而是把限制条件直接嵌入裁决句：例如先承认四次
多项式的 $R^2$ 达到 1，紧接着写明样本只有五个点，因而不把插值结果等同于整体拟合；
又在切换方法前保留旧方案“结果不理想”的具体对象。写作时可复用“虽然……，但由于
……，故选取……”“在原模型的基础上加入……，调整后 $R^2$ 由……变为……”和“考虑
到上述规划结果……，下面改用……”等推进关系，但省略处必须填本题读数、限制和动作。
B160 进一步给出更细的比较和诊断承接：“当……相同、……相同，只改变……时，将……
组进行对照”先证明组间可比；“由图……可以看出……；在同一温度下……；从整体来看……”
把两个响应分开读取；“从回归的相对效果看……；从绝对效果看……；从方差分析方面看……”
则按指标功能安排句序。这些表达之所以自然，是因为每个分句都填有保持量、组名、读数或
统计量，而不是因为用了更多连接词。

本轮补入的十五篇又把“人味”落到了模型接口上。A070、A147、A195、A212 的共同点是
先把机理候选、控制体内核、分层目标或基线残差写清，再说明哪一项证据触发降阶、增量、
第二层裁决或局部修正；A028、A115 则把近共线量级和硬约束超限直接写成几何简化或目标
改写的理由。B078、B108、B125、B175 不从“综合决策”起笔，而是先区分信息集、在线/离线
状态、库存接口和玩家目标，只有信息边界改变后才引入随机转移或 MDP。C109、C142、C170、
C227、C305 的段落都保留字段缺口和模块交接：评级或违约标签缺失时先补输入，风险排序
再进入额度、利率和放贷决策，疫情只改影响层。可调用的短句是“问题二的目标没有变化，
缺口只在……；先补……，再把输出接回……”，以及“前问离线路径不能直接复用，因为当前
状态只保留……”。这些句式的前提是本题确实存在相同接口和已点名的缺口，不能用来装饰
没有跨问依赖的段落。

B030 的惯用承接更适合几何分类和长仿真结果。“根据……将相对位置关系分为如下三类”
先保证覆盖；“但在求解过程中，发现……在几何上具有对称性，可通过旋转重叠”再给压缩
依据；“固定……，遍历……，记录所得全部结果，并绘图输出”把程序动作改写为可核的论文
步骤。需要复用时，作者先写“由问题一（1）中的定位模型4可知”或“联合式……得”，
而不是直接说“同理可得”。这些词句只在分句中填入具体形态、变量、范围、式号和输出时
才有人味；留下省略号的空壳仍是模板腔。

B035 又补出一组更接近日常推导的转折词。作者用“事实上”“不难发现”“注意到”“同理”
承接几何事实，但每次后面都紧跟圆、弦、交点或角度区间；需要换方法时，先说原来四区域
分类为何麻烦，再写“这里我们应用一个相对更简单的求解方法”。误差论证用“仍旧以……
为例”“下面我们通过……进行说明”把例子、图形和区间表逐步接起。结果比较不回避局部
反例，而以“在局部上可能……占优势，但从全局上看……”转入误差较大点先被修正的机制。
这些惯用词可用，但只有补入当前题目的对象、数值和下一步动作时才不同于 AI 式空转。

B086 的惯用承接更适合处理信息边界、算法比较和几何流程。作者先用“该问本质上是一个
……问题”给小节定性，再以“考虑到……故需……”把限制接到动作；局部搜索处明确写
“由于……信息不共享，故无法通过全局……”，随后用“注意，此处定义的目标函数不是
全局函数”阻止概念混用。定义事后评价量时又补一句“这是合理的，因为……并没有根据
全局误差函数做出调整”，把评价和决策分开。预处理证明后使用“由此可见……更重要的
是……”从误差界转到搜索次数；方案裁决前使用“特别注意……无实际参考价值，其仅用于
比较……”限定资源代理。稳健性段以“需要注意的是，单凭一组原始数据并不能准确检验
……故……”引出十组扰动。这些表达的人味来自它们公开了限制、用途和下一步，而不是
来自句首词本身。

B226 的惯用承接集中在“先归约、再命名、后推导”。“由问题一知，只要已知……和……，
就可以求解……，因此我们需要……”把复用接口写得很清楚；“为了让模型的建立过程更加
直观易懂，我们记……”让面、线和向量先拥有共同称谓；“根据……规程……出于建模的
严谨性，我们在下文证明……”把经验规范和数学授权分开。贪心求解用“以……为始边”“覆盖
宽度恰好能达到……”“以 1 m 为步长”逐项写出程序动作；求解器切换则先说三个未知参数
及其较大范围会使循环遍历耗时，再引出粒子群。分区段主动使用“我们人为地将……”并列出
依据，拟合不佳时又以均方误差触发局部再分。结果段采用“区域 V 的测线极多。这是因为……
又因为……这就需要……”的区域特定因果链。这些句式只有填入旧模型、两个缺失量、几何
对象、边界、步长、误差指标和异常区域时才可复用；若只留下“因此”“这是因为”，仍是 AI
式模板。

B311 的自然感首先来自它不掩去定义上的犹豫。作者先承认“由题意可以得到两种理解”，
分别推到可计算结果，再用共同数值例和 CAD 图中的可观察差异作裁决；“模型一”“模型二”
不是并列堆砌，而是一真一假的候选。跨问处的“比较本问与问题一可知”后面则紧跟不变
结构、改变的水深和夹角，以及需要替换的函数输入。小坡度也没有被一句“影响较小”略去，
而是先换算到 2 海里，看到累计差值后再决定方向。方向选定以后才补一句与测量规范一致，
把题内判断和外部印证分开。

本篇还说明，有序连接词并非天然带有 AI 腔。“首先、其次、再次、最后”在原段中分别
检查对称、单调、特殊方向不变和极限式，四个分句的主语与裁决对象都不同，因此读者
能逐项核对。递推段用“从西边界开始”“当第 34 条……”“故第 35 条舍弃”完成对边闭合；
最不利位置段则先找等深线间距最大处，再写“若该处完全覆盖，其余位置亦可覆盖”。可
迁移的不是这些句首词，而是“条件--数值--边界--动作”均不省略的组织方式。

B477 的惯用承接常出现在作者发现原路走不通的时候。“……的结果是反直觉的，若继续
使用，……将失去意义”先把旧定义的失效对象说透，随后才写“为适应……，将……改为
……”。“由于前述流程尚未确定从浅处还是深处起算，两种情况均作计算”则公开未定选择，
再以同一越界面积裁决。方法切换处的“这种做法在后续计算中有以下几点负担”后面列的是
四个可核对象，而不是“效率较低、效果不佳”的同义反复；新方法紧接着给出深度、梯度、
坡度和覆宽四个下游量。路线段的“经过初步尝试，发现……”也保留起点依赖与漏测，继而
用“因此只对……”限定返工范围。可迁移的骨架是“具体现象--受影响对象--必要修改--
交付接口”，不能只摘取“反直觉”“针对上述问题”“经过尝试”这些句首词。

C006 的段落推进更接近真实业务建模。作者先承认“供货连续性”不能直接读取，再找间隔和
连续周数作代理；在周能力估计处不用历史均值，不是泛称“考虑周期性”，而是分别写出均值
过于均匀和作物具有周期两个理由。遇到“优化结果没有材料偏好”这一反常现象时，段落也没有
立即换算法，而是先检查成本尺度和等产能关系，再局部修改 A/C 约束。结果部分把 24 周总体
走势与两个具体周次并列，并用“最大约降 $5\%$、平均仅降 $0.956\%$”限定改进强度。这些句子
自然，是因为它们允许“不采用”“不明显”和“继续沿用”出现，并各自给出理由。

C085 的人类痕迹集中在“具体失败--局部修补”和“先改比较口径”两种动作。作者不写
“为提高算法性能引入动态权重”，而是先说全局固定权重会造成“总体情况良好但是某一周
极度缺少供货”，再按该周到期望产能的差更新下一状态权重；也不泛称“考虑周期性”，
而是点出第 7、8 周连续淡季，由此取两周前视。类似地，“历史成本略低”没有直接成为
优劣结论，因为两套方案的产能不同；换成成本/产能后才作裁决。这里可复用的惯用承接是
“本文设定了两个规则……在以上两个规则都满足的情况下……”“若全局只使用……可能会
导致……因此……”“经过数据分析，发现第……周……因此……”和“……虽然略低，但由于
……未达到……故取……作为判断标准”。每个省略处都必须填阈值、周次、基准和数值。

C169 的惯用词服务于“先给理由，再决定纳入或舍弃”。构造非对称指标时，作者连用“为了
体现出……的影响明显重于……”“由于……的正负造成的影响不同，所以……满足下列条件”
和“综合考虑以上条件和生活实际，构造出……”，顺序是业务差异、数学条件、具体函数。
发现供货规律后又写“容易发现……但在下文的分析中不单独分析，原因有如下两点”，允许
分析以舍弃结束。初始状态用“考虑到企业并非刚刚成立……因此不能简单地认为……为 0”，
可分性用“由于……无记忆……所以……；为了简便计算可以一周为单位考虑”。这些表达像人写，
不在于连接词本身，而在于每个连接词后都跟着可核对的业务事实和建模后果。

C283 的自然感更多来自作者把犹豫、否决和边界留在正文中。假设段不是直接宣布简化成立，
而是写“由题可知……；若假设……，则……难以满足……，故原假设不成立”；权重段用“但根据
实际操作结果来看……与事实不吻合，需要在适度范围内修正”承接客观结果与业务判断。算法段
先分别写 LP 和 DP 的具体代价，再落到双层贪心；达到库存边界后用“见好就收”收束，并紧接
继续增加订单的损失。随机结果用“多次实验”“频率较高且目标较优的典型代表”限定结论强度；
方向权重用“一方面……不可过小；另一方面……不可过大，否则……”同时写出下界和上界。
这些惯用词不能脱离对象照抄，真正应复用的是“提出判断--暴露反例或代价--限定范围--作出选择”。

C155 的自然感来自作者愿意写出“为什么这里不做”。对颜色缺失，正文不是机械插值，而是先说明
相同已知条件仍出现多种颜色，再指出随机表面取样不能代表整件文物，随后落到“因此本文认为不能
通过已知数据填补……，后文涉及颜色时忽略……”；对表单 3，又用“与表单 2 不同的是……”说明
风化字段已经可直接观察，因而删除重复估计步骤。卡方段常写“由表……可知……满足使用 Pearson
卡方检验前提，因此……；……不满足……，因此……”，把检验名称放在前提之后。回归段以“在后续
预测过程中，仅对 $R^2$ 较高的成分进行预测，对较低成分采取不变策略”把指标直接接到动作；异常
距离段则按“当回归方程确定时，距离 $d$ 是唯一决定变量--异常值会导致预测失真--为了减小影响而
压缩距离”推进。可复用的不是句首“因此”，而是每个因此之前都有字段、样品编号、阈值或失真机制，
之后都有剔除、保留、换检验、预测或不预测的具体动作。

C229 的自然感则来自作者把同一字段内部的证据差异写出来。颜色填补不是先定方法再套数据，而是
先写“按此分为两个类别分别进行填补”，逐一列出 40 对 39、58 对 11 的参照关系；19、48 号
因候选化学样品“参考性较低”，又改用众数。近似零值段用“不能单纯将缺失值插补为 0，可看作
近似零值”承接检测限和舍入语义；作图段直接承认“因二氧化硅占比较大，图中不能很好地体现其余
成分”，故保留全图并补画去除图。亚类段先写“没有已知因变量，因此该类问题为无监督学习问题”，
结果不清时又保留“仅可粗略将其分为三类”。可复用的是“候选质量触发方法分流--图形局限触发
补图--有无标签决定方法身份--证据强弱决定语气”，不是热卡、聚类或“故”字本身。

C050 的惯用承接更接近真实的数据分析过程。正文先写“考虑到数据较为庞大和复杂，为更好地分析
和理解数据”，随后不是泛泛宣布预处理，而是逐项观察丰富度、销量、损耗和分布；对象过多时又写
“由于……种类较多，本文选取……”并立即交代代表规则和附件去向。回归段的“由回归系数可知”
后面紧接 $1\,\mathrm{kg}$ 与价格变化的业务读法，特征段的“为有效判断……，本文选取……”后面
则给出对象、口径和代理含义。这些句子有人味，是因为作者在每次缩减数据、选择对象或解释系数时
都公开了判断依据，而不是因为“考虑到”“由于”“由此可知”出现得多。

更值得保留的是作者没有把代理量写成唯一解释。打折次数既可能对应受欢迎和供不应求，也可能对应
销售较弱或保鲜期较短；文章先并列竞争机制，再要求补看销量、反馈与竞品信息，最后分别落到补货、
定价或促销动作。现有附件与待采外部数据也分别成段：前者支持已经量化的关联，后者只说明采集方式、
可辨机制和决策用途。可迁移的段落不是“还需采集更多数据”这一空句，而是“现有证据能回答什么--
哪两种机制仍无法区分--补哪项观测--不同结果分别采取什么动作”的完整判断轨迹。

C126 的惯用承接更像作者在一边处理数据、一边向读者说明为什么这样处理。正文先用“由于……，
直接采用……会……”保留不正态、零值和计算负担，再写“为避开……同时降低……，用同期平均法
将……”；“为此确定”后面跟的是新时间序列和迭代设置，不是重复结论。对象太多时，作者常写
“这里仍以……为例进行说明”，并在下一句交代大矩阵为何不宜展开；参数表之前又用“因篇幅限制，
这里只给出……”明确正文保留什么、支撑材料保存什么。这些句式自然，是因为每个省略处都要填
数据现象、计算代价、代表对象和实际文件，而不是把“由于”“为此”“因篇幅限制”当成润色词。

结果段的承接更值得模仿。“由实际经营过程可知”后面紧跟隔日难售这一事实，并据此确定 VAR(2)；
“表中第一行……第二行……第三行……”先教读者读取系数、标准差与检验量，才综合 AIC、SC、
$R^2$ 和 F 值把结论限定到未来一周；“注：负数表示……”则把负预测转成余货、折价和不补货。
最后，“从另外一个方面进行筛选”不是平行堆模型，而是用半小时销售频率独立核对 29 个对象。
这组表达真正可复用的，是“事实--尺度--选型”“表项--指标--限定结论”“数值--状态--动作”和
“主路线--独立规则--有限一致”四条判断轨迹。

C228 的惯用词同样依附于具体判断。开放题先写“由于题目未明确给出具体的探究方向，需要先考虑
从哪几个角度入手”，随后立即列出时间、销售和关系三个观察面；这里的“由于”不是装饰，而是在
解释为何正文先聚焦而不是直接建模。预测段用“首先采用……，结果发现效果较差，因此……”保留
候选方法的失败状态；单品段则写“观察最近筛选出的……，发现……”后逐项指出缺失、无明显趋势
和随机波动，再把复杂模型降为加权平均。两段都让“结果发现”后面跟可观察事实，让“因此”后面
跟可执行选择。

求解器切换时，作者先写“与问题二不同”，再列出本问新增的 $0/1$ 进货标志、连续决策和双目标，
最后才交代混合编码与 NSGA-II；效果解释则先说明为什么选早七月的历史经营区间，再把方案利润放入
这个基线中。改进建议用“尽管……直观便捷，但其主要从……出发”承认成本导向定价的作用，再从
其边界推出竞品、预期价格和收入数据。应迁移的是“开放任务--观察面”“候选效果--换路”“数据
不足--降级”“变量结构--求解器”“参照理由--有限比较”和“现有方法--边界--补采”的段落次序，
不是机械复制“由于、首先、因此、尽管”。

A0127 的自然感来自作者愿意在算法之前交代工程判断。开头用“它实际上是一道优化问题”给全题定性，用“问题二是问题一的反问题”说明问间关系；随后不是笼统说计算困难，而是先把题面公式中可直接得到的量剔除，只留下五类效率。塔影段的“考虑到”后面有明确负担、数量关系和中心对称，因而“以镜面中心作判断”是一个有来源的近似；布场段的“从图中可以看出”后面紧接北侧高效、镜面北移和塔位南移三个动作；问题三的“若直接采用遍历法……存在计算量过大的问题”又自然引出先二分、再变步长。灵敏度段还保留“快速增加至一定值后几乎不变”的饱和描述。这里应模仿的是“事实—负担—取舍—动作—边界”的信息顺序，而不是复制几个句首。

A0165 的惯用表达更像作者在顺着计算过程说话。摘要常先写“将……划分为……，分别进行建模”，随后用“通过……确定……”“根据……计算……”“利用……判断……”逐级交付中间量；问题分析则从“为计算……首先需要……”起笔，先补齐太阳位置和镜面法向，再进入遮挡、截断和功率。作者写简化时不只说“为提高效率”，而是交代“分析其中一个变量时，将其他变量固定”，并用完整计算作基准说明 $2\times2$ 栅格的偏差。图形段的“由第一问散点图可知”后面紧跟塔向南移动这一设计动作；问题三的“在问题二的基础上”后面则列出不再改变的宽度、塔位和位置，只把重心放在各圈高度。检验段先说明前两问由第一问变化而来，再改变采样时刻核对峰值移动是否符合太阳视运动。应模仿的是每个常用词后都有具体对象、冻结量、动作和裁决，不把“通过、根据、因此”留下空壳。

A016 常用“由于……因此应当满足……”“同时由于相邻……只需要……，因此选取……”“若……则撤回上一步，并将步长除以 2”“但是由于我们按照时间遍历……立即停止，因此……”推进。这里连续出现的“由于、因此、但是”并不显得模板化，因为每个词都在承担不同责任：第一个引入物理顺序，第二个裁决多根，第三个规定数值回退，第四个用事件先后排除后发分支。“可以发现……唯一确定，因此……固定不变”则用来从几何关系推出没有可优化自由度；“不妨设……变为 $k$ 倍，由……可知……也变为 $k$ 倍”把局部比例传到全链。真正该写进 Skill 的不是这些词的表面频次，而是连接词后必须出现条件、被选对象、执行动作和结论边界。

可迁移的段落骨架目前归纳为一百五十类。第一类是观察段：数据或图形中出现了什么现象，该现象让
哪个假设或方法变得合理。第二类是改进段：原模型遗漏什么结构，这一遗漏会造成何种
偏差，因而补入哪个变量、交互项或约束。第三类是比较段：固定数据、指标和预算，
列出候选结果，再说明采用其中一个的依据。第四类是实验设计段：先写预算和安全
边界，再列候选方法、舍弃理由、因素水平与最终组合。第五类是成对结果段：在同一
对象下先给基准或无约束结果，再给受限或调整结果，最后解释差值对决策的含义。
第六类是小样本裁决段：先报看似占优的指标，再交代样本数、自由度或外推限制，最后
说明为何保留较低阶方案。第七类是递进指标段：复述上一轮指标和剩余不足，写出本轮
新增结构，再用新指标决定继续、停止或换用另一类方法。第八类是双响应控制变量段：
先写保持条件、唯一改变量和比较组，再分别读两项响应，只作局部判断。第九类是诊断
修模段：初始模型先点名多重共线性等具体问题，给出逐步回归等修正动作，随后按相对
拟合、绝对误差和整体检验三个口径解释修正结果。第十类是潜变量代理段：先说明业务概念
不可直接观测，再列代理指标、方向和权重解释。第十一类是替代量否决段：先写方便的候选
口径，再用两个不重叠的理由说明它会丢失什么信息，最后给替代构造。第十二类是反常结果
修正段：结果与预期不符时先诊断机制，只修改受影响的约束或系数。第十三类是状态传递段：
明确本期输入、决策、实现量和期末状态如何形成下一期初值。第十四类是宏观--微观解释段：
总体序列回答方案是否闭合，代表时点说明状态和约束如何具体作用。第十五类是有限效果段：
同时给最大值与平均值，直接承认改善幅度有限，并解释可行域或资源边界。第十六类是模型沿用
段：逐项核对输入、输出和目标是否改变，再决定复用、修改、删除或新增模块。
第十七类是规则筛样段：先说明每条阈值分别排除何种不稳定，再规定同时满足后的统计量。
第十八类是失败驱动改权重段：点名固定权重造成的局部缺货状态，再由当前缺口更新下一轮权重。
第十九类是数据定窗段：用连续淡季或峰谷位置确定前视/后视长度，并解释库存调配方向。
第二十类是同口径裁决段：当成本、收益或误差对应的产出规模不同时，先改成单位产出指标再比较。
第二十一类是非对称损失段：先写两类偏差的业务后果差异，再列函数必须满足的方向与曲率条件，
最后构造分段表达式。第二十二类是观察后舍弃段：先说明规律和可能成因，再用干扰来源与任务变化
解释为何不纳模。第二十三类是初值与可分性段：先判断系统是否已有历史状态、过程是否有时间记忆，
再决定初值来源和能否逐期拆分。第二十四类是保守聚合段：先明确局部解的业务语义，再说明跨场景
应取最大还是最小，以及该方向怎样保证统一可行。第二十五类是反证假设段：先提出便利假设，
再代入容量、数量或需求事实得到矛盾，最后修正假设。第二十六类是权重异常修正段：先给客观权重，
点名与业务不符的指标，再限定主观修正幅度并做一致性检查。第二十七类是贪心停止段：先定义满意
边界，达到即停，并说明继续选择会增加的采购、运输或库存代价。第二十八类是候选机制裁决段：
所有候选共享历史基准，分别说明变动机制，再按约束、目标、初期稳健性和现实含义逐个淘汰。
第二十九类是随机代表段：报告重复运行，只以频率和目标双条件选典型方案，不宣称结果唯一。
第三十类是目标逆转扩围段：先比较新旧目标，再解释先前低排名对象为何重新有贡献，据此调整候选范围。
第三十一类是分类压缩段：先列完整情形，再指出可旋转、平移或交换重合的等价类，最后给保留分支和
适用条件。第三十二类是计算动作段：按固定对象、遍历对象、范围/步长、记录项和图表产物交代仿真，
使代码循环在正文中有清楚语义。第三十三类是长结果分流段：说明完整结果的存放位置、正文选取的代表
情形和代表理由，避免以“篇幅有限”掩盖证据。第三十四类是结构转化复用段：先完成坐标、图结构或
子问题分组的转化，逐项核对旧模型接口，再指出哪些推导不重述、哪些特殊对象仍需另建关系。
第三十五类是负担触发换模段：先保留旧方法及其分类或计算负担，再指出仍可利用的不变量，最后
给出新构造和被消去的工作量。第三十六类是数值到区间证明段：由具体例子定位可能区域，再检查
轨迹端点、给出变量界，并用互不交叠的区间或表格闭合穷举。第三十七类是局部让步--全局反转段：
先承认某算法在一轮或一个局部指标上占优，再说明状态更新或优先处理大误差点怎样改变总体目标。
第三十八类是任务充分停止段：同时列数值上还能达到的精度、任务阈值、额外轮次或资源代价，说明
为何停止在满足任务而非机器精度的位置。第三十九类是复用前失效审计段：先说明新问本质上复用
旧模型，再集中列出共线、共圆或接口退化，逐项给检测条件和修复动作。第四十类是信息边界账本段：
先写对象之间不共享什么，再列局部可观测量和局部目标；若需要全局指标，只把它作为事后评价，并
解释为何没有进入对象决策。第四十一类是预处理双收益段：先证明误差界收紧，再说明搜索区域、函数
评估或信号次数怎样减少。第四十二类是双轴或三轴方案裁决段：在同一初始数据下依次比较收敛速度、
稳定误差和资源代理，公开代理口径后再选择。第四十三类是引理到流程装配段：逐条写固定点、调整点、
控制角/比例和几何目标，再映射到具体对象与阶段。第四十四类是单样本不足段：先承认一组初值不足以
检验稳定性，再定义扰动分布、重复数据集、收敛区间和结论范围。第四十五类是缺失语义裁决段：
先判断空值在题面中表示未检出、未知还是未采样，再用同条件多值性和采样范围决定置零、拒绝填补或
局部忽略。第四十六类是检验前提分流段：先列候选检验的期望计数等前提统计量，再按满足与否选择
Pearson、Yates 或其他方法。第四十七类是双向证据段：用前问趋势和已知标签解释后问聚类，再用
聚类结果反向核对前问结论。第四十八类是指标门控计算段：逐对象报告 $R^2$ 等指标，把高低分支
分别接到预测、保持不变或改用基线。第四十九类是异常距离局部修正段：先点名唯一敏感量和异常放大
造成的失真，再只对该量作单调压缩，不借此替换整条主模型。第五十类是变量--样本双聚类段：R 型
聚类先降相关并选代表变量，Q 型聚类后分样本，分别说明两个对象和输出。第五十一类是数据语义删步
段：比较新旧表单的采样或标签语义，确认新表已直接给出状态后，删除重复估计而复用后续模型。
第五十二类是同基准类型比较段：固定一个参考序列，分别报告各类最大、最小关联及集中/离散程度，
最后解释类型差异。第五十三类是同字段分组填补段：先列匹配条件和候选样品，再按参考质量选择
热卡或众数，并写出具体编号与填充值。第五十四类是近似零值辨析段：区分空白和报表零的来源，
说明为何不能作普通零，再接乘法替换与闭合。第五十五类是主导成分遮蔽拆图段：保留全量图，
另画去除主导成分的图，结合两图读取大小量变化。第五十六类是可解释规则--多变量复核段：
简单阈值先给直观分类和测试结果，PLS-DA/VIP 再检查多变量差异与一致性。第五十七类是无监督
身份判定段：没有已知标签或因变量时，先说明不能监督分类，再用曲线选择聚类数。第五十八类是
聚类强弱诚实限定段：并列两类分簇图，一类清楚便具体描述，另一类不清楚只作粗分并限制结论。
第五十九类是双指标定维段：同时比较交叉验证误差与累计解释率，选择最小充分维数。第六十类是
对称矩阵配对检验段：先由对称性取上三角去重，再由成分对对应关系选择配对检验并报告决策。
第六十一类是分问关系定性段：先交代研究对象与本质身份，再依次说明正向计算、反问题和放宽约束。
第六十二类是关键困难隔离段：从题面公式中剔除可直接获得的量，只把剩余未知量交给后续计算模块。
第六十三类是对称与统计补偿简化段：先写精确边界的负担，再由数量占优、中心对称或统计补偿推出
中心判据和二值动作。第六十四类是坐标转移判定段：对象由 A 坐标转至公共坐标，再转入 B 坐标，
通过平面交点和区域内外完成有效/无效判断。第六十五类是能量收支闭合段：从全反射能量扣除无效点
损失，再以接收能量或射线数形成效率比值。第六十六类是空间证据到布场段：由云图或方向性读出
高效区域，把镜面集中方向、塔址反向移动、特征距离和分区层规则连续写出。第六十七类是计算预算到
粗细搜索段：先说明直接高精度遍历为何过重，再由二分或大步长定位初步范围，缩短步长精化并给停止
精度。第六十八类是范围--方向--饱和--排序段：明确扰动对象与步长，报告变化区间、增减非对称、
饱和位置和分指标敏感度次序。第六十九类是过程分段摘要：把物理或业务链拆成连续阶段，每一阶段
使用一个动作动词并交付一个下游量。第七十类是等价几何准备段：先给分解关系，再把原判定改写为
可计算的子判定，最后说明二者为何等价。第七十一类是近邻栅格仿真段：先限定候选近邻，再写射线、
平面和区域判定，避免把全场两两检查塞入一句。第七十二类是快速核误差门段：以高精度计算为基准，
说明低分辨率核节省什么，并用同一输入报告偏差后限定其用途。第七十三类是固定其余变量降维段：
列出固定量，只让一个变量改变；从响应曲线确认形态后，再选择一维搜索。第七十四类是空间图触发
方向变量段：由散点或云图读出高效方向，把方向转成设施移动、坐标符号和搜索区间。第七十五类是
继承旧解并开放新自由度段：先说明前问已经确定的变量，再逐项冻结，只把新增自由度交给本问求解。
第七十六类是结构预言检验段：先由机理写出图形或峰值应怎样变化，主动改变输入结构，重新计算后
以观察到的结构与预言是否一致作裁决。第七十七类是条件分支传播段：先写第一阻断条件，成立时
明确跳过的后续计算；不成立再进入第二、第三判定，只有全部通过才计算接收或输出。第七十八类是
邻接降算量段：先给全对全对象数、网格与角度共同形成的规模，再由物理局部性设距离阈值并只遍历
候选集合。第七十九类是抽样粗细寻优段：小比例多次抽样负责估计粗参数，冻结弱变量后提高比例、
减小步长并承担最终评价。第八十类是弱影响变量冻结段：固定基线分别试算变量，报告指标变化区间，
取代表或较优值后把搜索预算转向关键变量。第八十一类是趋稳阈值选步长段：列出分辨率序列和指标，
指出增幅开始变缓的位置，用“接近稳定”限定结论，再给工作步长。第八十二类是达标结果反推目标段：
方案超过硬阈值后，分析删去低效或高干扰对象对总量、面积、干扰和单位量的连锁作用，重新决定
次目标并验算可行性。
第八十三类是唯一剩余量隔离段：把题面数据代入既有模型，逐项列出可直接量，只留下唯一不能直接
求得的量，解释随机或高维原因后再给抽样法与样本数。第八十四类是重叠区域离散段：先指出两区域
可能重叠、直接相加会重复计数，再给离散对象、网格步长和并集计数规则。第八十五类是空间猜想
复核段：由序号或时间曲线提出位置猜想，另选代表日期和时刻绘制二维分布，复核后才转成布局动作。
第八十六类是方向抵消降维段：固定其余参数沿两个方向分别扰动，依据对称性、跨时段抵消和影响
幅度冻结弱方向坐标，只保留主方向。第八十七类是前问结构按组降维段：前问得到规则结构，本问
先说明沿用，再由对称性让同组对象共享参数，把逐对象变量压成逐组变量并重新选择求解器。
第八十八类是可行样本背景检验段：在相同硬约束下生成 $N$ 组可行基线，以同指标、同单位比较，
只陈述选定方案在该样本背景中的优势，不外推为全局最优。
第八十九类是两个关键量归约段：先写旧模型的输入输出接口，点名两个缺失几何量，分别推导后
代回旧模型，不用“沿用前问”省略接口。第九十类是命名对象后推导段：先命名平面、直线和向量，
再写位置/垂直关系、叉乘、坐标表达和物理意义。第九十一类是规范加自证段：规范说明常用做法，
命题把最优性压成同长度候选比较，覆盖面积或同口径指标承担裁决。第九十二类是边界锚定贪心段：
选起始边界，以恰达边界确定首个对象，按固定空间步长保留可行候选，取最低可行重叠率并迭代至
对边。第九十三类是绝对--相对双误差检验段：基线与随机对照比较同一位置或指标，同时报告绝对
误差和相对误差，只作同口径一致性结论。第九十四类是复杂度触发换求解器段：先列未知量个数、
取值范围和遍历成本，再给群智能算法、参数和停止规则。第九十五类是残差驱动区域细化段：按结构
初分区域，拟合局部代理并读取 MSE，只拆分失败区、重拟合子区而保留合格区。第九十六类是分区域
结果解释段：点名异常区域和指标，用局部坡度、覆盖宽度或作业简化解释线数、间距和漏测后果。
第九十七类是竞争定义裁决段：承认题意歧义，保留两个候选定义，固定共同数值例，寻找可区分的
观察量，再写保留与舍弃。第九十八类是改变量复用段：比较前后两问，列不变结构和改变量，只修改
受影响的函数输入。第九十九类是小参数长尺度累积段：把局部小量换算到完整任务尺度，以累计差值
和两端可行区间判断能否忽略。第一百类是内部裁决加外部印证段：先由题内数值比较作选择，再说明
结果与标准或规范一致。第一百零一类是四项结构检查段：依次核对对称、单调、特殊方向不变和极限式，
每一项对应独立预测。第一百零二类是对边闭合递推段：从首边界和首个对象开始，按固定规则递推，
计算最后可行对象的对边余量，并舍弃越界候选。第一百零三类是读者连续性插图段：后文确需调用
早先几何时，说明为减少回翻而就近重现，只保留当前推导所需标记。第一百零四类是最不利位置全局
覆盖段：定位等值线间距最宽或最平缓处，说明其最不利性，由该处充分条件推出全域覆盖并核边界。
第一百零五类是侧写到决策接口段：逐项写侧写对象、观察、所支持的变量/窗口/约束和后续动作，
不让探索性图表与模型脱节。第一百零六类是代表规则先行段：对象过多时先给排序指标、阈值和代表
选取规则，再展示代表结果并标出全量附件位置。第一百零七类是系数业务量纲段：从回归系数读出
自变量每变化一个业务单位时因变量的方向和变化量，再说明经营含义。第一百零八类是特征语义账本段：
依次写特征对象、计算口径、代理现象和下游用途，不以特征名称代替构造理由。第一百零九类是现有--
外部证据分流段：先列附件内能够量化的判断，再列尚未覆盖的信息、采集方式和预定用途。第一百一十类
是代理竞争机制段：同一代理量保留至少两个可行解释，指出能够区分它们的追加证据，并分别给出动作。
第一百一十一类是主文--附件规则分工段：正文保留选择规则、代表结果、关键数字和解释，附件承接
全量表、补充图及代码，且正文明确给出去向。第一百一十二类是数据病态触发换口径段：依次写
分布、零值、时间结构和计算负担，再给转换规则、新数据口径与下游模型。第一百一十三类是判断矩阵
压缩段：先定义阈值、显著性和字符串含义，再以 T/F 等接口代替正文逐格复述，并给全量系数去向。
第一百一十四类是业务时间尺度授权段：从保鲜、余货或运营时滞得到记忆长度，由此确定阶数或窗口，
不以“综合考虑”跳过尺度推导。第一百一十五类是检验表导读段：先说明各行列对应估计量、不确定性
和检验量，再按信息准则、拟合和显著性分别裁决，最后限定任务与时窗。第一百一十六类是异常数值
业务翻译段：把负值、零值或极端量先映射为库存、退货、折价等可观察状态，再决定保留、截断或转成
动作。第一百一十七类是低成本独立核对段：主模型给出候选后，用频率、阈值或业务规则独立筛选，
比较名单和差异，只在口径可比时称有限一致。第一百一十八类是预测--优化中间量接口段：逐项写
预测量进入目标或约束的位置，并把近期均值、上期状态、损耗、成本和业务下限接到下一阶段决策。
第一百一十九类是开放题聚焦段：先说明题面没有指定探究方向，再把对象拆为能够互相补足的观察面，
并交代这些观察面共同覆盖什么任务。第一百二十类是候选失败换模段：列出实际试过的同类方法、当前
指标或现象上的不足，再说明新方法正好处理哪项困难。第一百二十一类是数据能力触发降级段：写清
缺失、短窗、无趋势或随机波动为何使复杂模型不可靠，再给简单估计及其时窗边界。第一百二十二类是
变量结构授权求解器段：比较前后问的连续、离散和多目标结构，只在变量域或目标关系改变后切换算法。
第一百二十三类是弹性修正接口段：先给基线预测，再用响应或弹性函数修正，逐项指出修正量、成本、
损耗和补货量进入目标或约束的位置。第一百二十四类是同期历史基线段：先解释为何选同季、同星期或
同业务状态，再把方案数值放入该参照范围作有限比较。第一百二十五类是从模型边界生长建议段：先写
当前模型由哪些内部信息驱动，再指出它不能区分的现实关系，最后列能区分关系的数据和动作。第一百
二十六类是扰动范围限定段：写明被扰动输入、范围、其余保持量和观测指标，只把结论限定为当前扰动
与报告范围内的稳定性检查。
第一百二十七类是局部递推展开段：先求起点，再定义从第 $i$ 个对象到第 $i+1$ 个对象的关系，最后说明如何沿空间与时间铺开。第一百二十八类是物理条件选根段：先列顺序、方向或单调条件过滤数学根，再以局部连续性选最近可行根。第一百二十九类是首次事件缩域段：由结构重复、轨迹继承或单调关系限定最早事件涉及的对象，并说明为何不会漏掉更早情形。第一百三十类是事件顺序消枝段：算法遇到首次事件即停止时，明确哪些只会随后发生的分支无需进入当前判定。第一百三十一类是越界回退求精段：写清初始步长、前进判据、越界回退、缩步比例、停止阈值和由此得到的误差界。第一百三十二类是固定量否决优化段：先由相切、守恒或代数约束证明目标量唯一，说明当前约束下没有优化自由度，再进入指定方案。第一百三十三类是跨区域复用段：复用前问公式时逐项列出不等式、坐标符号、切向量、初值和边界的改变。第一百三十四类是齐次传递缩放段：先证明局部状态改变 $k$ 倍会传到相邻状态，再用已有解与上限求整体比例。第一百三十五类是同输出灵敏度比较段：对多个输入分别作同尺度扰动，始终观察同一输出，最后按变化方向和幅度给敏感度排序。
第一百三十六类是反证逐节前推段：先假设第 $i$ 个构件最先触发事件，再由轨迹滞后推到
$i-1$，逐节前推至首构件并否定假设。第一百三十七类是事件对象切换解释段：参数微变而
输出跃变时，先核对首次触发对象是否改变，再用局部结构差异解释不连续。第一百三十八类是
终点可行反例段：对只检查终态的简化方案，沿全过程搜索更早事件，以一个明确反例裁决其
适用性。第一百三十九类是现实动作推出边界段：说明参数继续变化会要求哪个构件执行怎样的
不可能动作，再据此给出上下界。第一百四十类是异常时段局部解释段：先圈定异常时段，再比较
同一时间内不同节点经过的路程，由局部轨迹解释速度或响应变化。

第一百四十一类是立场先于检验段：先写决策者更在意哪类错误，再确定原假设、备择假设、功效
和样本量。第一百四十二类是无样本先交代模拟器段：公开缺失的数据、模拟分布和用途，只把
模拟结果限定为流程比较。第一百四十三类是结构规模授权搜索段：列出节点、边或方案数量，
说明精确枚举负担后再切换搜索方法。第一百四十四类是下游后果触发预处理段：不先说“去噪”，
而先说明干扰会怎样改变反演或决策，再给处理规则。第一百四十五类是有限改善收窄结论段：
主量基本不变而局部指标改善时，只称局部修正，不写“全面提升”。

第一百四十六类是共同环境复评段：各候选先在各自情景中生成，再放入同一随机环境比较，避免
把情景差异误当方案差异。第一百四十七类是管理语言参数化段：把“优先、稳定、兼顾”等口号
翻译为可调参数或优先级，并说明参数改变哪条约束。第一百四十八类是安全--及时--可执行损失段：
先分别定义三种错误后果，再选择阈值或动态规划目标。第一百四十九类是样本行转主体段：重复
测量资料先聚合到真实决策主体，再计算首次达标时点，避免行数替代人数。第一百五十类是总体--
少数类--样本数并列段：分类结果同时给总体指标、关键少数类召回和对应样本量，不用单一准确率
结束段落。

句子节奏也有可辨认的规律。定义和假设用短句，避免一个句子同时承担五个动作；
推导段允许稍长，但每个公式后立即解释变量或几何意义；结果段先给条件和数字，
下一句再说明趋势或机制。图表不以“如下图所示”结束，而写成“由图可知，温度升高
后……先……后……”，或按“图左/图右”分别完成比较。相比之下，“构建综合性框架”
“充分证明模型有效”“显著提升了性能”“为后续研究奠定基础”等 AI 惯用语若没有
具体对象、数值和证据，一律不用。

上述结论由 59/59 篇论文的文风入口支持：其中四十四篇保留完成原页核验的细粒度
句式和段落节奏，十五篇补入结构化判断动作卡。它们的作用是提供可调用的写作动作，
不把人工选出的高信息段落冒充 59 篇全文词频。每条惯用表达均与观察、限制、动作或结果绑定；
若省略处填不进本题事实，便不调用该句式。

\section{创新切入、高频失分与高分策略}

\subsection{从基线到差异化模型}

创新先回答“基础方案在哪个可验证环节失败”，再引入复杂度。可执行路径包括：
\begin{enumerate}
  \item \textbf{结构创新：}把题意中的时空、网络、分层或阶段结构显式写入状态和约束；
  \item \textbf{机理—数据融合：}用守恒方程约束参数范围，以数据校准未知项，并分开报告拟合和外推；
  \item \textbf{基线—改进链：}先给解析、贪心、线性或持续性基线，再针对残差、非线性或不确定性改进；
  \item \textbf{约束驱动：}将现实可行性写入模型，并报告活跃约束、不可行诊断和修复代价；
  \item \textbf{不确定性驱动：}把参数误差传播到结果和决策，寻找稳定区间与策略切换点；
  \item \textbf{预测—决策耦合：}把预测概率及其校准误差传入目标和约束，比较 oracle 标签、预测标签和概率输入下的策略变化；
  \item \textbf{验证创新：}增加独立回代、网格收敛、消融、滚动验证或反事实情景，使结论比模型名更可信。
\end{enumerate}

不采用“换一个更新颖算法名称”作为独立创新。若鲸鱼、麻雀或蛾火焰优化相对
确定性基线、GA/PSO 没有同预算优势，就不进入主模型；若深度网络没有超过季节
或树模型基线，保留简单方案；若多指标组合权重没有通过扰动和秩稳定检验，不把
组合本身写成贡献。
\subsection{高频失分点}

本节不再逐篇搜集疏漏，也不虚构出现频数。以下八项来自 59 篇优秀论文中反复需要闭合的写作
接口，同时也是竞赛中一旦缺失便会削弱说服力的质量门。检查时只判断自己的稿件是否闭合。

\begin{enumerate}
  \item \textbf{题问没有闭合。}摘要、问题分析、模型、结果表和附件应逐问对应；每问至少留下
  “对象--方法--结果--解释”四个可核对位置。
  \item \textbf{选择没有来由。}写出模型或算法前，先交代观察到的结构、上一方法的具体不足、
  可行域规模或误差目标；删掉“综合考虑”“为了提高精度”后仍应看得懂为何换法。
  \item \textbf{数学对象没有落到题意。}变量、索引、单位、边界和约束都应能翻译回真实对象；
  公式后用一句话说明各项作用，结果后报告活跃约束或余量。
  \item \textbf{求解过程不可复现。}正文至少给初始化、参数、停止规则和输出；随机算法还要给
  预算、种子或重复统计，精确算法给界、间隙或可手算小例。
  \item \textbf{结果只有数字没有判断。}按“条件--读数--比较--原因--行动”解释图表；遇到跳变、
  异常时段或反直觉排序，先查事件对象、局部结构和口径变化，不用“结果合理”收尾。
  \item \textbf{检验与主模型同源。}机理题用守恒、极限、网格收敛或独立几何核对；优化题用
  约束回算、基线和共同情景复评；数据题按实体或时间隔离验证，并同时报告关键少数类。
  \item \textbf{不确定性没有传到决策。}参数扰动、预测概率、情景分布和中间分类误差应进入
  目标或约束；区分固定策略复算与重新优化，并给稳定区间或策略切换点。
  \item \textbf{正文与附件断开。}表图、代码和数据产物使用同一版本标识；正文给入口和关键
  参数，附录给依赖、数据字典、运行顺序与完整输出，修改结果后重新生成表图和摘要数字。
\end{enumerate}

这些门的用途是约束自己的稿件，不是给 59 篇论文排错。真正应模仿的仍是作者如何把观察写成
判断、把判断写成动作，再用结果和检验收住这一动作。
\subsection{高分写作策略与冻结节奏}

高分策略可落实为五个冻结点：定题后冻结任务矩阵；基线通过手算/小规模枚举后
冻结变量和口径；主模型通过独立检验后冻结关键数字；图表与附件可追溯后冻结
结果；最后才冻结摘要。任何改动关键结果的代码都必须重新生成正文表图，避免
人工复制导致不一致。评阅者首先需要看见的是问题闭合、模型合理、数字可核验和
结果有解释，而非模型数量。

提交前执行双人交叉核对：一人只看题问—摘要—结果是否对应，另一人只看假设—
变量—公式—代码—附件是否对应；最后由未撰写该段的人检查引用、单位、有效
数字、图表位置和页面溢出。审计脚本用于发现结构和引用问题，但不能替代公式
正确性、学术判断和原始页面核验。

\section{三项工具成果及迭代机制}

\subsection{学习报告}

本文件 \texttt{assignment.tex} 使用 \texttt{ctexart} 编写，含完整导言区、
标题、正文、公式、长表、参考材料和附录，可由 XeLaTeX 独立编译。报告只陈述
模板和 Skill 的完成情况，不把两者全文嵌入，三项成果保持独立。

\subsection{双模式论文模板}

模板实际路径为 \path{thesis template/}。\par
目录保留上游
\path{cumcmthesis.cls}、字体和
\path{example.tex}；\path{main.tex} 默认为指导模式，在每级标题后给出目标、
清单、语言、篇幅排版、图表位置、判别力要点和错误警示；\texttt{submission.tex}
定义模式开关后载入同一主文件，隐藏全部指引。章节顺序固定为摘要与关键词、
问题重述、问题分析、模型假设、符号说明、分问题模型建立与求解、结果分析与
模型检验、灵敏度/稳健性、模型评价与改进、参考文献、附录。

篇幅提示集中使用本轮 59 篇论文的实测中位数和四分位范围，不预造比例。目录
\texttt{ITERATION.md} 记录证据范围和改造日期；后续范文结论先在批注区登记
来源编号、日期、页码、题型和核验状态，再增量修改指引。

\subsection{MCM-CUP-STANDARD-WRITE Skill}

Skill 位于
\texttt{C:/Users/Lenovo/.codex/skills/mcm-cup-standard-write/}，调用名为
\texttt{\$mcm-cup-standard-write}，界面名保留
\texttt{MCM-CUP-STANDARD-WRITE}。核心文件写入任务规定的两条第 0 条及第 1--7
条约束，详细资料拆为论文结构、A/B/C 题型、章节句式、模型创新、质量门和证据
更新六类主题 reference，并另设人类文风、语料索引与逐篇证据卡。当前五十九张正式卡已经把
十四类章节、模型链、结果与检验、创新、高分机制、惯用词/段落骨架和必要原图边界逐篇录入。公开审计接口为：
\begin{quote}
\small\ttfamily
python audit\_manuscript.py MAIN.TEX\\
\hspace*{1em}--problem-type A\textbar B\textbar C\\
\hspace*{1em}--format text\textbar json
\end{quote}
脚本检查章节缺失和顺序、占位符、图表标签、交叉引用、文献引用、结果解释、
模型检验及附录结构。它的输出是质量门线索，不能证明模型结论正确。

\section{结论}

以“算法清单”为起点，模型、代码和写作链才是本轮可验收的产物。三名成员分别
承担机理连续模型、综合优化和数据决策的首责，编程、误差分析和选题判断则共同承担；
18 道真题分别提供练习入口。论文研究以 59 篇编号语料为限，59 张可追溯的正式证据卡覆盖
2892 页；赛题和专家评述不混入文风样本，A/B/C 三类分别完成 20/19/20 篇人工门禁。
模板与 Skill 可随本轮 59 篇证据增量修改；四周计划把尚无实证的能力落实为代码、图表、
论文和复现产物，并逐项验收。后续限时演练不增加算法名称，重点检验基线建立、约束书写、
关键结果复算和文字说明能否在时限内完成。

\begin{thebibliography}{99}
\bibitem{papers}
2020--2025 年全国大学生数学建模竞赛 A/B/C 题编号论文语料（本地归档），
59 篇，2892 页。

\bibitem{problems}
全国大学生数学建模竞赛赛题原文：2020--2025 年 A/B/C 题（本地归档）。

\bibitem{commentaries}
全国大学生数学建模竞赛专家评述：2020--2025 年相关题目（本地归档）。

\bibitem{mindmap}
数学建模学习思维导图 \texttt{1.png}（用户提供，本地文件，核验日期：
2026-07-28）。
\end{thebibliography}

\appendix
\section{算法责任与验收总表}

\begin{longtable}{p{2.1cm}p{5.0cm}p{2.0cm}p{5.5cm}}
\caption{思维导图算法的完整分配}\label{tab:algorithm-owner}\\
\toprule
模块 & 算法或模型 & 主责 & 最小验收\\
\midrule
\endfirsthead
\toprule
模块 & 算法或模型 & 主责 & 最小验收\\
\midrule
\endhead
缺失与异常 & 均值/中位数、插值、KNN、删除；$3\sigma$、箱线图、Isolation Forest；重复样本删除/合并 & 组员2 & 处理前后审计表、规则理由、下游敏感性\\
尺度与编码 & Z-score、Min--Max、one-hot、label、等宽/等频分箱 & 组员2 & 训练折拟合、逆变换或字段映射测试\\
探索与相关 & Pearson、Spearman、Kendall、分布和关联可视化 & 组员2 & 假设条件、显著性/效应量和非因果声明\\
聚类 & K-Means、DBSCAN、层次聚类、谱聚类 & 组员2 & 合成样例、簇稳定性和业务解释\\
插值与降维 & Lagrange、三次样条、Kriging；PCA、t-SNE、UMAP & 组员2 & 留点误差、解释方差、可视化边界\\
样本筛选 & Monte Carlo 随机抽样 & 组员2（与组员1共享接口） & 固定种子、抽样分布、覆盖率和偏差审计\\
力学与物理 & Newton、弹簧阻尼、刚体平衡、摩擦、重力、Ohm/Kirchhoff、Bernoulli/质量守恒、理想气体 & 队长 & 受力/系统图、量纲、守恒/回代检查\\
ODE 模型 & Logistic、捕食—被捕食、Newton 冷却、药物代谢、SIR/SIS & 队长 & 初值、参数、解析/小样例和积分误差\\
PDE 模型 & 热、波、扩散、多孔流 & 队长 & 初边值、离散格式、稳定条件和网格收敛\\
数值方法 & Euler、RK4、FDM、FEM & 队长 & 同题对照、误差阶或收敛证据\\
规划 & LP、NLP、QP、整数、MILP、0--1、单/多目标、确定/随机模型 & 组员1 & 小规模枚举、界/间隙、约束映射\\
智能优化 & GA、PSO、SA、ACO、DE、贪心、NSGA-II、MOPSO、WOA、SSA、MFO & 组员1 & 同预算基线、20 种子、可行率和分布\\
图与网络 & Dijkstra、Floyd、SPFA、Prim、Kruskal、最大流、最小费用流、路径规划、二分图匹配、节点重要性 & 组员1 & 标准图测试、复杂度和负边/容量边界\\
仿真与敏感性 & Monte Carlo、情景仿真、参数扰动与策略切换 & 组员1 & 采样依据、重复次数、置信区间和稳定区间\\
主观评价 & AHP、模糊综合评价、FAHP & 组员2 & 判断矩阵、一致性、隶属函数与权重扰动\\
客观评价 & 灰色关联、PCA、因子分析、熵权、CRITIC、变异系数 & 组员2 & 正向化、尺度、权重来源和秩稳定性\\
综合评价 & TOPSIS、VIKOR、灰色综合评价、物元可拓、熵权--AHP、CRITIC--TOPSIS & 组员2 & 方法语义、基线排序和扰动对照\\
回归 & 多元线性、Ridge、Lasso、SVR、RF、XGBoost、CatBoost & 组员2 & pipeline、交叉验证、残差与外推边界\\
分类 & Logistic、朴素 Bayes、SVM、决策树、RF、AdaBoost、CNN & 组员2 & 混淆矩阵、阈值、类别不平衡与校准\\
时间序列 & ARIMA、SARIMA、指数平滑、VAR、GM(1,1)、灰色--Markov、长期预测、Prophet、XGBoost、LightGBM、BP、RNN、LSTM、TCN、GRU & 组员2 & 持续/季节基线、滚动验证、预测区间和外推边界\\
验证与误差 & MAE、MSE、RMSE、MAPE、K 折、留出、训练测试、分组/滚动验证 & 三人交叉 & 划分先于预处理；指标与任务损失一致\\
工程实现 & MATLAB、Python、配置、日志、随机种子、图表和结果导出 & 三人交叉 & 一键入口、依赖文件、单元测试和结果追溯\\
\bottomrule
\end{longtable}

\section{证据文件说明}

全量 OCR 页面记录、论文汇总、证据台账、模型链和语料统计属于临时分析证据，
保存在工作区 \texttt{.cumcm-work/} 下；正式成果不依赖运行时缓存才能编译。
台账中的页面文本可用于定位，不替代原 PDF。任何涉及公式、表格、小字号图注
或低置信页的结论，须回到原始页面局部核验后才能加入模板或 Skill。

\end{document}
